% zeta2_7.tex -- "The 2-adic zeta value zeta_2(7) is irrational".
% Build: pdflatex zeta2_7 (twice). Self-contained: no \input files, the
% bibliography is contained in this file.
\documentclass{amsart}
\usepackage{amsmath,amssymb,amsthm}
\usepackage{alltt}
\usepackage{array}
\newcolumntype{P}[1]{>{\raggedright\arraybackslash}p{#1}}
\usepackage[hidelinks]{hyperref}
\hypersetup{
  pdftitle={The 2-adic zeta value zeta2(7) is irrational},
  pdfauthor={Gennaro Marco Devincenzis}}

\newtheorem{theorem}{Theorem}[section]
\newtheorem{lemma}[theorem]{Lemma}
\newtheorem{proposition}[theorem]{Proposition}
\newtheorem{corollary}[theorem]{Corollary}
\newtheorem*{mainthm}{Main Theorem}
\newcommand{\thisthmname}{}
\newtheorem*{genericthm}{\thisthmname}
\newenvironment{namedthm}[1]{\renewcommand{\thisthmname}{#1}\begin{genericthm}}{\end{genericthm}}
\theoremstyle{definition}
\newtheorem{definition}[theorem]{Definition}
\theoremstyle{remark}
\newtheorem{remark}[theorem]{Remark}

\numberwithin{equation}{section}
\setcounter{tocdepth}{1}
\setlength{\emergencystretch}{1.5em}

\newcommand{\Z}{\mathbb{Z}}
\newcommand{\Q}{\mathbb{Q}}
\newcommand{\N}{\mathbb{N}}
\newcommand{\R}{\mathbb{R}}
\newcommand{\sgn}{\operatorname{sgn}}
\newcommand{\Res}{\operatorname{Res}}
\newcommand{\Lev}{\operatorname{Lev}}
\newcommand{\Top}{\operatorname{Top}}
\newcommand{\dd}{\,\mathrm{d}}
\newcommand{\half}{\tfrac12}

\begin{document}

\title[The $2$-adic zeta value $\zeta_2(7)$ is irrational]{The $2$-adic zeta
value $\zeta_2(7)$ is irrational}

\author{Gennaro Marco Devincenzis}
\urladdr{https://github.com/gmDevi}

\thanks{\emph{Disclosure of AI use.} The mathematics of this paper (the
construction and the proofs), the Lean formalization (apart from the
vendored third-party files for the prime number theorem \cite{PNTLean}) and
the text were produced by AI agents (Anthropic Claude models) under the
direction of the author. The proofs were checked by adversarial AI
verification and, in their formal version, by the Lean kernel. No human expert
has refereed them. The formal theorem statement is machine-checked; the
identification of the $2$-adic zeta value used in Lean with the
Kubota--Leopoldt value is cited from the literature and is not formalized.
Details are in Section~\ref{sec:ai}.}

\date{28 September 2026}

\subjclass[2020]{Primary 11J72; Secondary 11S40, 11S80, 11M06, 68V20}

\keywords{$p$-adic zeta values, irrationality, Volkenborn integral,
well-poised rational functions, denominators, formal verification}

\begin{abstract}
Let $\zeta_2(s)=L_2(s,\omega^{1-s})$ be the $2$-adic zeta values of Kubota and
Leopoldt. We prove that $\zeta_2(7)$ is irrational. The proof follows the
approach of Lai, Sprang and Zudilin for $\zeta_2(5)$: Volkenborn integrals of
the third derivative of explicit rational functions $R_n$ give linear forms
$S_n=\rho_n+Z_n\zeta_2(7)$ with rational coefficients, and Lai's $2$-adic
irrationality criterion concludes. The rational function $R_n$ has six nested
well-poised numerator blocks and pole orders that ramp up from $1$ to $4$
towards the middle, and it is integrated at the centred point
$t+\frac12-\frac n2$; as a result, no zeta value other than $\zeta_2(7)$
occurs. The $2$-adic size of $S_n$ is bounded with the $\Delta$-calculus of
Sprang and Lai, the size of the coefficients by a residue-level estimate, and
the nonvanishing follows from an auxiliary prime $\ell=\frac{19}{20}n-1$, after
Lai and Sprang. The main new ingredient is a bound for the odd part of the
common denominator of $\rho_n$ and $Z_n$, obtained from a cancellation among
the residues in each residue class modulo a prime $q$ (a class lemma together
with a global residue identity). Two steps of the proof are computer-assisted:
the maximum of an explicit potential, certified by interval and exact
rational arithmetic, and the value of the denominator exponent, an exact
rational computation over finitely many intervals. The theorem is formalized
in Lean~4 with Mathlib, without hypotheses and with only the standard axioms;
the one input that is cited rather than formalized is the identification, due
to Lai, Sprang and Zudilin, of the Volkenborn-integral definition of
$\zeta_2(7)$ used in Lean with the Kubota--Leopoldt value. The mathematics,
the proofs and the formalization were produced by AI agents under the
direction of the author, and have not been refereed by a human expert.
\end{abstract}

\maketitle

\tableofcontents

% ===========================================================================
\section{Introduction}\label{sec:intro}

\subsection{Background}\label{ss:background}

For a prime $p$ and an integer $s\ge2$ let
\[
  \zeta_p(s):=L_p(s,\omega^{1-s}),
\]
where $L_p(s,\chi)$ is the Kubota--Leopoldt $p$-adic $L$-function and $\omega$
is the Teichm\"uller character (\S\ref{ss:zeta2}). The number
$\zeta_p(s)\in\Q_p$ is the $p$-adic limit of $\zeta(k)$ as $k\to s$
$p$-adically through negative integers $k\equiv s\pmod{p-1}$, and
$\zeta_p(s)=0$ for even $s$. The odd values $\zeta_p(3),\zeta_p(5),\dots$ are
$p$-adic analogues of the odd values of the Riemann zeta function.

In the archimedean case, $\zeta(3)$ is irrational by Ap\'ery's theorem
\cite{Ape79}, infinitely many of the values $\zeta(2k+1)$ are irrational by
the theorem of Rivoal and Ball--Rivoal \cite{Riv00,BR01}, and at least one of
$\zeta(5)$, $\zeta(7)$, $\zeta(9)$, $\zeta(11)$ is irrational by a theorem of
Zudilin \cite{Zud01}; the irrationality of $\zeta(5)$ is open. On the $2$-adic
side, Calegari \cite{Cal05} proved that $\zeta_2(3)$ and $\zeta_3(3)$ are
irrational, and Beukers \cite{Beu08} gave another proof. Calegari, Dimitrov
and Tang proved that $\zeta_2(5)$ is irrational with their arithmetic
holonomy method \cite{CDT20,CDT25}. Lai, Sprang and Zudilin \cite{LSZ26} gave a
second proof by an explicit construction in the spirit of Ap\'ery, which also
bounds the irrationality measure of $\zeta_2(5)$.

Sprang \cite{Spr20} used Volkenborn integrals of rational functions to
construct linear forms in $1$ and $p$-adic zeta values, and proved a $p$-adic
analogue of the Ball--Rivoal theorem. Lai and Sprang \cite{LS25} showed that
many of the numbers $\zeta_p(3),\zeta_p(5),\dots,\zeta_p(s)$ are irrational
(at least a constant times $\sqrt{s/\log s}$ of them, for large odd $s$). The
$2$-adic analogue of Zudilin's theorem is due to Lai \cite[Theorem~1.1]{Lai25}:
for every integer $s\ge0$, at least one of the numbers $\zeta_2(j)$ with $j$
odd and $s+3\le j\le3s+5$ is irrational; for $s=3$ this says that at least
one of $\zeta_2(7)$, $\zeta_2(9)$, $\zeta_2(11)$, $\zeta_2(13)$ is irrational. Lai, Lupu and Sprang
\cite{LLS25} proved analogous statements for all primes $p\ge5$. The author's
Lean formalizations \cite{Dev26,DevPair26} contain proofs that at least one of
$\zeta_2(7)$, $\zeta_2(9)$, $\zeta_2(11)$ is irrational and that at least one
of $\zeta_2(7)$, $\zeta_2(9)$ is irrational; both were produced by AI agents
and are registered in the Palomar registry of Lean-verified results
\cite{Pal26,PalPair26}.

\subsection{The result}\label{ss:result}

We call an element of $\Q_2$ \emph{rational} if it lies in the image of
$\Q\to\Q_2$.

\begin{mainthm}
The $2$-adic zeta value $\zeta_2(7)$ is irrational.
\end{mainthm}

The normalization of $\zeta_2(7)$ plays no role: the values of Coleman
\cite{Col82} and Furusho \cite{Fur04} differ from $\zeta_2(7)$ by the non-zero
rational factor $(1-2^{-7})^{-1}$ (\S\ref{ss:zeta2}). The Main Theorem implies
the two statements of \cite{Dev26,DevPair26} quoted above, and the cases
$1\le s\le4$ of Lai's theorem (those whose range contains $7$).

An unrefereed draft by C.~D.~Long \cite{Lon26a}, posted on GitHub in August
2026, claims that $\zeta_2(s)$ is irrational for every odd $s$ with
$3\le s\le29$, which includes $\zeta_2(7)$, by a different method (a hybrid
form of arithmetic holonomy). Our proof does not use that draft. We were
unable to verify it: an internal review, carried out with AI assistance and
published with the formalization (\texttt{docs/long\_review.md}), found that its
large-prime estimate (the torsor decomposition and the proposition on ``torsor
Cartier strictness'') does not appear to hold for the Eisenstein--Eichler leaf,
and that without it the method's margin for $\zeta_2(7)$ is negative; the
draft's other ingredients and all of its printed numerical margins were
reproduced. The review is not a refereed report, and the author has not been
contacted. No refereed proof of the irrationality of $\zeta_2(7)$ is known to
us (\S\ref{ss:novelty}).

\subsection{Outline of the proof}\label{ss:outline}

Fix $n\in40\Z_{>0}$, let $\eta=(\eta_1,\dots,\eta_6)$ with
$40\eta=(5,7,7,9,11,13)$, and let
$a_k\in\{1,2,3,4\}$ ($0\le k\le n$) be the pole orders that equal $1$,
$2$, $3$ at distance less than $\frac n{20}$, $\frac n{10}$, $\frac{3n}{20}$
from the ends $\{0,n\}$ and $4$ in the middle (\S\ref{ss:construction}). Put
\begin{equation}\label{eq:intro-R}
  \begin{gathered}
  R_n(t)=(2t+n)\,\frac{\prod_{m=1}^{6}\bigl(t+\frac12+\eta_mn\bigr)_{(1-2\eta_m)n}}
  {\prod_{k=0}^{n}(t+k)^{a_k}},\\
  S_n=\int_{\Z_2}R_n'''\Bigl(t+\frac12-\frac n2\Bigr)\dd t,
  \end{gathered}
\end{equation}
where $(x)_L$ is the rising factorial and the integral is the Volkenborn
integral. The rational function $R_n$ has degree $-3$ and satisfies
$R_n(-n-t)=-R_n(t)$; its poles have order at most $4$. These three properties
imply that $S_n=\rho_n+Z_n\zeta_2(7)$ with $\rho_n,Z_n\in\Q$, and that no
other zeta value occurs (Lemma~\ref{lem:L}). The proof then consists of four
estimates, which are combined by Lai's criterion in the form of Lai and
Sprang (\S\ref{sec:main}).
\begin{enumerate}
\item[(V)] \emph{$2$-adic size} (Theorem~V, \S\ref{sec:V}):
  $v_2(S_n)\ge6.8n+7-10(\lfloor\log_2\frac{3n}4\rfloor+1)$, by the
  $\Delta$-calculus of Sprang and Lai \cite{Spr20,Lai25}.
\item[(A)] \emph{Odd denominators} (Theorems~C, K and~A, \S\ref{sec:den}):
  the least odd $D_n$ with $D_n\rho_n,D_nZ_n\in\Z[\frac12]$ satisfies
  $\frac1n\log D_n\le E^*+o(1)$ with $E^*=3.8808431\ldots$
\item[(G)] \emph{Archimedean size} (Theorem~G, \S\ref{sec:G}):
  $\log\max(|\rho_n|,|Z_n|)\le U^*n+O(\log n)$ with $U^*=0.70670223$.
\item[(N)] \emph{Nonvanishing} (Theorem~N, \S\ref{sec:N}): if
  $\ell=\frac{19}{20}n-1$ is prime, then $v_\ell(\rho_n)=-1$ and
  $v_\ell(Z_n)\ge0$.
\end{enumerate}
The criterion needs $E^*+U^*<6.8\log2=4.7134008$; the margin is
\[
  E^*+U^*-6.8\log2=-0.1258555\ldots\quad\text{per }n.
\]

The construction was found by a numerical search among rational functions of
this shape; among the features that make it work, three seem essential. First,
the \emph{centred evaluation} at $t+\frac12-\frac n2$: the Volkenborn
translation formula then produces correction terms at half-integers $h$ with
$|2h|\le n-1$ only, and this is what allows us to show that no prime $q>n$
divides the denominators (Theorem~C), without any extra zeros of $R_n$ at
half-integers.
Secondly, the \emph{pole-order ramps}: the orders $1,2,3,4$ near the ends
lower both the growth and the denominators, and orders at most $4$ leave only
$\zeta_2(5)$ and $\zeta_2(7)$ as candidates among the odd zeta values, of
which $\zeta_2(5)$ drops out because $R_n$ has degree $-3$. Thirdly, the \emph{nested well-poised blocks}: six blocks
centred at $\frac{n-1}2$, of lengths $\frac34n$, $\frac{13}{20}n$ (twice),
$\frac{11}{20}n$, $\frac9{20}n$, $\frac7{20}n$, whose total length $3.4n$
balances the total pole order $3.4n+4$.

The main new ingredient is the denominator bound (A). Classical arguments
bound the $q$-adic valuation of each residue of $R_n$ separately; this gives
the ``residue level'' $\max(0,4-E_\xi)$, where $E_\xi$ is the balance of zeros
and poles in the residue class $\xi$ modulo $q$, and it leads to the
exponent $5.08$, far above the break-even value $4.0067$. We expand the
contribution of each residue class modulo $q$ in a local variable and show
that it equals, up to higher-order terms, a residue of $R_n$ times an explicit
kernel (the \emph{class lemma}, Lemma~\ref{lem:class}). The sum of these
residues over all classes vanishes (Lemma~\ref{lem:kernel}), because $R_n$
has degree $-3$. Subtracting it gains up to several powers of $q$ in each
class at once. The resulting per-prime bound (Theorem~K) is a finite
computation on the local models of the classes, which depends only on
$n/q$ and $k_0/q$; a continuum version of it is piecewise constant on
explicit intervals (Lemmas~\ref{lem:piecewise}--\ref{lem:periodic}), and the
prime number theorem turns it into the exponent $E^*$ (Theorem~A).

\subsection{Computer assistance}\label{ss:intro-computer}

The informal proof is computer-assisted at two places (\S\ref{sec:computer}):
the maximum $U^*$ of an explicit potential on $[0,1]$ (Lemma~\ref{lem:G7}),
certified with interval arithmetic and with exact rational arithmetic, and the
value of $E^*$ (Theorem~A), an exact rational computation of a step function
on $1449$ explicit intervals plus a periodic tail. Both were redone by
independent implementations. Only the inequality $E^*+U^*<6.8\log2$ matters,
and it holds with a margin of $2.7\%$ of the $2$-adic gain. In the formal
proof all numerical inputs are checked by the Lean kernel
(\S\ref{sec:formal}).

\subsection{Formal verification}\label{ss:intro-formal}

The Main Theorem is formalized in Lean~4 \cite{MU21} with Mathlib
\cite{Mathlib20}. The Lean theorem \texttt{Zeta2.Z7.zeta2\_7\_irrational} has
no hypotheses and depends only on the three standard axioms
\texttt{propext}, \texttt{Classical.choice} and \texttt{Quot.sound}. The Lean
development defines $\zeta_2(7)$ through Volkenborn integrals as $J_6/768$,
with the same trusted definitions as the registered formalization
\cite{DevPair26,PalPair26}; the identification of this number with
$L_2(7,\omega^{-6})$ is \cite[Lemma~2.8]{LSZ26}, which is cited and not
formalized. The formal proof follows the text with slightly weaker constants
(growth rate $0.71$, denominator rate $3.99$). Section~\ref{sec:formal}
gives the statement, the trusted definitions and the trust base.

\subsection{Novelty and related work}\label{ss:novelty}

A literature search by AI agents on 24 and 26 September 2026 (arXiv, including
the papers citing \cite{Lai25} and the author pages of Lai, Sprang and
Zudilin; OpenAlex; Zenodo; GitHub; public talks and blogs; a general web
search for ``$\zeta_2(7)$'') found no refereed paper and no arXiv preprint that
proves the irrationality of $\zeta_2(7)$. The only public claim found is the
unrefereed draft \cite{Lon26a} mentioned in \S\ref{ss:result}, which uses a
different method. Google Scholar, MathSciNet, the full text of zbMATH and
Lai's thesis were not searched, so we cannot rule out other work.

The methods are those of \cite{Spr20,Lai25,LS25,LSZ26}: Volkenborn integrals
of rational functions, the $\Delta$-calculus, Lai's criterion and the
auxiliary-prime condition of Lai and Sprang \cite[Lemma~2.2]{LS25}.
Non-uniform pole and zero structures and ``inserted'' rational zeros have
archimedean ancestors in the work of Zudilin \cite{Zud04} and of Fischler,
Sprang and Zudilin \cite{FSZ19}, and well-poised symmetry goes back to Ball
and Rivoal \cite{BR01} and Rhin and Viola \cite{RV01}. We did not find prior
$p$-adic uses of the following ingredients, which we therefore believe to be
new: the centred evaluation point; the pole-order ramps together with nested
well-poised blocks giving a form in $1$ and a single zeta value; the
residue-class cancellation for the denominators (Lemmas
\ref{lem:class}--\ref{lem:farpart} and Theorem~\ref{thm:collective}); and the
auxiliary prime $\ell=\frac{19}{20}n-1$. This is a statement about what we
found, not a proof of absence.

\subsection{Organization}\label{ss:organization}

Section~\ref{sec:prelim} collects definitions and published tools.
Section~\ref{sec:construction} gives the construction and the linear form.
Sections~\ref{sec:V}, \ref{sec:den}, \ref{sec:G} and~\ref{sec:N} prove the
four estimates, and Section~\ref{sec:main} proves the Main Theorem.
Section~\ref{sec:computer} describes the computer-assisted steps and their
certification, Section~\ref{sec:formal} the formalization, and
Section~\ref{sec:ai} how this work was produced with AI agents and what has
been verified, and how.

% ===========================================================================
\section{Preliminaries}\label{sec:prelim}

\subsection{Notation}\label{ss:notation}

Throughout, $q$ and $\ell$ denote primes. We write $v_q$ for the $q$-adic
valuation on $\Q$ and on $\Q_q$, with $v_q(0)=+\infty$, $|x|_q=q^{-v_q(x)}$,
and $\Z_{(q)}=\{x\in\Q:v_q(x)\ge0\}$. The archimedean absolute value is
$|\cdot|$ and $\log$ is the natural logarithm.
\begin{itemize}
\item $(x)_0=1$ and $(x)_k=x(x+1)\cdots(x+k-1)$ (rising factorial);
  $\binom{x}{k}=x(x-1)\cdots(x-k+1)/k!$.
\item $s_2(N)$ is the binary digit sum of $N\ge0$; Legendre's formula gives
  $v_2(N!)=N-s_2(N)$, and $v_q(m!)=\sum_{e\ge1}\lfloor m/q^e\rfloor$.
\item $[\varepsilon^\lambda]F$ is the coefficient of $\varepsilon^\lambda$ in a
  power series or Laurent series $F(\varepsilon)$, and $[y^p]$ likewise.
\item $\theta(x)=\sum_{q\le x}\log q$ is Chebyshev's function.
\item For an integer $k\ge1$, $k_-$ is $k$ with its leading binary digit
  deleted, so $k-k_-=2^{\lfloor\log_2k\rfloor}$.
\end{itemize}
We use the following elementary facts: Kummer's bound
$v_q\binom NK\le\lfloor\log_qN\rfloor$; among $L$ consecutive odd integers, at
least $\lfloor L/q^e\rfloor$ are divisible by $q^e$ ($q$ odd); and for an odd
prime $q$, $\sum_{u=1}^{q-1}u^{-s}\equiv0\pmod q$ unless $(q-1)\mid s$.

\subsection{The Volkenborn integral}\label{ss:volkenborn}

We follow \cite[Ch.~5]{Rob00}, \cite[\S2.2]{Lai25} and \cite[\S2.2]{LSZ26}.
A function $f\colon\Z_2\to\Q_2$ is \emph{Volkenborn integrable} \cite{Vol72} if
the Riemann sums
\begin{equation}\label{eq:riemann}
  \frac1{2^M}\sum_{x=0}^{2^M-1}f(x)
\end{equation}
converge $2$-adically as $M\to\infty$; the limit is
$\int_{\Z_2}f(t)\dd t$. The integral is $\Q_2$-linear but not translation
invariant. Strictly differentiable functions are integrable \cite[p.~264]{Rob00};
in particular rational functions without poles in $\Z_2$ are.

\begin{lemma}[Translation; {\cite[Lemma~2.2]{Lai25}, \cite[Lemma~2.4]{LSZ26}}]\label{lem:translation}
Let $f$ be strictly differentiable on $\Z_2$ and $d\ge1$ an integer. Then
\[
  \int_{\Z_2}f(t+d)\dd t=\int_{\Z_2}f(t)\dd t+\sum_{l=0}^{d-1}f'(l),
  \qquad
  \int_{\Z_2}f(t-d)\dd t=\int_{\Z_2}f(t)\dd t-\sum_{l=-d}^{-1}f'(l).
\]
\end{lemma}

The second formula follows from the first applied to $t\mapsto f(t-d)$. For
$s\ge1$ put
\begin{equation}\label{eq:Js}
  J_s:=\int_{\Z_2}\frac{\dd t}{(t+\frac12)^s};
\end{equation}
the integral exists because $-\frac12\notin\Z_2$.

\subsection{\texorpdfstring{$2$}{2}-adic zeta values}\label{ss:zeta2}

We recall \cite[\S2.3]{LSZ26} (see also \cite[\S2.3]{Lai25} and
\cite[Ch.~11]{Coh07}). Let $\omega$ be the Teichm\"uller character and
$\langle x\rangle=x/\omega(x)$. For $s\ne1$ and $x\in\Q_2$ with $|x|_2\ge4$, the
$2$-adic Hurwitz zeta function is
$\zeta_2(s,x)=\frac1{s-1}\int_{\Z_2}\langle t+x\rangle^{1-s}\dd t$, and the
Kubota--Leopoldt $L$-function of a Dirichlet character $\chi$ of conductor
dividing $M$, $4\mid M$, is
$L_2(s,\chi)=\frac{\langle M\rangle^{1-s}}{M}\sum_{0\le j<M,\,2\nmid j}\chi(j)\zeta_2(s,j/M)$.

\begin{definition}[{\cite[Definition~2.7]{LSZ26}}]\label{def:zeta2}
For an integer $s\ge2$, $\zeta_2(s):=L_2(s,\omega^{1-s})$.
\end{definition}

This is the normalization of \cite{Lai25,LS25,LSZ26}. The values of Coleman
\cite{Col82} and Furusho \cite{Fur04} differ from it by the Euler factor
$(1-2^{-s})^{-1}$ \cite[\S2.3]{LSZ26}, a non-zero rational number. By
\cite[Lemma~2.4]{Cal05}, as quoted in \cite[\S2.3]{LSZ26}, $\zeta_2(s)=0$ for
even $s\ge2$.

\begin{lemma}[{\cite[Lemma~2.8]{LSZ26}}]\label{lem:LSZ28}
For every integer $s\ge1$, $J_s=s\,2^{s+1}\zeta_2(s+1)$. In particular
$J_s=0$ for odd $s$, and $J_6=768\,\zeta_2(7)$.
\end{lemma}

\subsection{The \texorpdfstring{$\Delta$}{Delta}-calculus}\label{ss:delta}

The following device of Sprang \cite{Spr20} and Lai \cite[Def.~2.3,
Lemmas~2.4--2.5]{Lai25} (see also \cite[Lemmas~2.5--2.6]{LSZ26}) bounds the
$2$-adic size of Volkenborn integrals. For $f\colon\Z_2\to\Q_2$ put
\[
  \Delta_0(f):=\inf_{k\ge1}v_2\Bigl(\frac{f(k)-f(k_-)}{k-k_-}\Bigr),\qquad
  \Delta(f):=\min\{1+v_2(f(0)),\Delta_0(f)\}.
\]

\begin{lemma}\label{lem:delta}
\begin{enumerate}
\item[(D1)] If $f$ is strictly differentiable, then
  $v_2\bigl(\int_{\Z_2}f\bigr)\ge\Delta(f)-1$.
\item[(D2)] If $f,g$ are $\Z_2$-valued, then $\Delta(fg)\ge\min(\Delta f,\Delta g)$;
  and $\Delta(f+g)\ge\min(\Delta f,\Delta g)$, $\Delta(Cf)=\Delta(f)+v_2(C)$.
\item[(D3)] A power series $\sum_ja_jt^j$ with $a_j\in\Z_2$, $a_j\to0$, has
  $\Delta\ge0$.
\item[(D4)] For $x\in\Z$ and $L\ge1$, $\Delta\bigl(\binom{t+x}{L}\bigr)\ge-\lfloor\log_2L\rfloor$.
\item[(D5)] $\bigl(\frac{\mathrm d}{\mathrm dt}\bigr)^\mu\binom{t+x}{L}=\sum_{r=\mu}^{L}c_{\mu,r}\binom{t+x}{L-r}$
  with $c_{\mu,r}=[z^r](\log(1+z))^\mu$ and $v_2(c_{\mu,r})\ge-\mu\lfloor\log_2r\rfloor$.
\end{enumerate}
\end{lemma}

\begin{proof}
(D1)--(D4) are \cite[Lemmas~2.4--2.5]{Lai25}; we recall the key identities.
(D1): $2^{-M}\sum_{k<2^M}f(k)=f(0)+\sum_{l<M}\sum_{k=2^l}^{2^{l+1}-1}(f(k)-f(k_-))/2^{l+1}$,
and each summand has valuation $\ge\Delta_0(f)-1$. (D2) is the product rule
for divided differences. (D4) follows from Vandermonde's identity
$\frac{f(k)-f(k_-)}{k-k_-}=\sum_{i=1}^{L}\frac1i\binom{k-k_--1}{i-1}\binom{k_-+x}{L-i}$
for $f=\binom{t+x}L$. (D5): $\frac{\mathrm d}{\mathrm dt}=\log(1+\nabla)$ on
the basis $\binom{t+x}{L}$, where $\nabla\binom{t+x}L=\binom{t+x}{L-1}$; the
coefficient $c_{\mu,r}$ is a sum of products of $\mu$ numbers $\pm1/r_i$ with
$r_i\le r$, times integers.
\end{proof}

\subsection{Irrationality criteria}\label{ss:criteria}

\begin{lemma}[{\cite[Lemma~2.1]{Lai25}}]\label{lem:lai}
Let $\alpha\in\Q_2$, and for $n$ in an infinite set $I\subseteq\N$ let
$a_n,b_n\in\Z$ and $L_n=a_n+b_n\alpha$. If
$\max(|a_n|,|b_n|)\cdot|L_n|_2\to0$ along $I$ and $L_n\ne0$ for all $n\in I$,
then $\alpha$ is irrational.
\end{lemma}

\begin{proof}
If $\alpha=u/b$ with $u\in\Z$, $b\ge1$, then $N_n=bL_n\in\Z\setminus\{0\}$,
so $1\le|N_n|\cdot|N_n|_2\le(b+|u|)\max(|a_n|,|b_n|)\,|L_n|_2$, a contradiction.
\end{proof}

The nonvanishing will come from an auxiliary prime, as in condition~(2) of
\cite[Lemma~2.2]{LS25}: if $\alpha=u/b$ and a prime $\ell\nmid 2b$ satisfies
$v_\ell(a_n)<v_\ell(b_n)$ (with $v_\ell(0)=\infty$), then
$v_\ell(bL_n)=v_\ell(a_n)<\infty$, so $L_n\ne0$. We write this out in
\S\ref{sec:main}.

\subsection{A valuation lemma}\label{ss:D2}

\begin{lemma}\label{lem:top}
Let $\Pi(\varepsilon)=\prod_\beta(1+\varepsilon/\beta)^{e_\beta}$ be a finite
product with $\beta\in\Q^\times$, $v_q(\beta)\ge0$ and $e_\beta\in\Z$. Then
$v_q([\varepsilon^\lambda]\Pi)\ge-\Top_\lambda$, where
\[
  \Top_\lambda:=\max\Bigl\{\sum_\beta j_\beta v_q(\beta):\ \sum_\beta j_\beta=\lambda,\
  j_\beta\ge0,\ j_\beta\le e_\beta\text{ whenever }e_\beta\ge0\Bigr\}.
\]
In particular $\Top_\lambda\le\lambda\max_\beta v_q(\beta)$; and if every
$\beta$ with $v_q(\beta)>0$ has $v_q(\beta)=1$ and $e_\beta\ge0$, and $s$ is
the number of such factors counted with multiplicity $e_\beta$, then
$\Top_\lambda\le\min(\lambda,s)$.
\end{lemma}

\begin{proof}
$[\varepsilon^\lambda]\Pi=\sum_{\sum j_\beta=\lambda}\prod_\beta\binom{e_\beta}{j_\beta}\beta^{-j_\beta}$
with $\binom{e}{j}\in\Z$, and $\binom ej=0$ for $0\le e<j$.
\end{proof}

\subsection{Prime number theorem}\label{ss:pnt}

We use the prime number theorem $\theta(x)=x+o(x)$, Chebyshev's bound
$\theta(x)\le1.02\,x$ for $x>0$ (which follows from
$\theta(x)<1.01624\,x$, \cite[Theorem~9]{RS62}), and Dirichlet's theorem for
primes $\equiv-1\pmod{38}$. The $o(x)$ is ineffective as used here.

% ===========================================================================
\section{The construction and the linear form}\label{sec:construction}

\subsection{The rational function}\label{ss:construction}

Fix $n\in40\Z_{>0}$.

\emph{Numerator blocks.} Let
$\eta=(\eta_1,\dots,\eta_6)=\frac1{40}(5,7,7,9,11,13)$ and, for
$m=1,\dots,6$,
\[
  C_m:=[\eta_mn,\,(1-\eta_m)n-1]\cap\Z .
\]
The block $C_m$ has $L_m=(1-2\eta_m)n$ elements, so the lengths are
$\frac34n$, $\frac{13}{20}n$, $\frac{13}{20}n$, $\frac{11}{20}n$,
$\frac9{20}n$, $\frac7{20}n$, with total $N:=\sum_mL_m=3.4\,n$. Each $C_m$ is
invariant under $c\mapsto n-1-c$, and every block contains
$[\frac{13}{40}n,\frac{27}{40}n-1]$. Write $\mu(c):=\#\{m:c\in C_m\}$.

\emph{Pole orders.} For $0\le k\le n$ put $\kappa_k:=\min(k,n-k)/n$ and
\[
  a_k:=\begin{cases}1,&\kappa_k<\frac1{20},\\ 2,&\frac1{20}\le\kappa_k<\frac1{10},\\
  3,&\frac1{10}\le\kappa_k<\frac3{20},\\ 4,&\kappa_k\ge\frac3{20}.\end{cases}
\]
Thus $a_k=1$ on $[0,\frac n{20}-1]\cup[\frac{19}{20}n+1,n]$, $a_k=2$ on
$[\frac n{20},\frac n{10}-1]\cup[\frac9{10}n+1,\frac{19}{20}n]$, $a_k=3$ on
$[\frac n{10},\frac{3}{20}n-1]\cup[\frac{17}{20}n+1,\frac9{10}n]$, and $a_k=4$
on $[\frac3{20}n,\frac{17}{20}n]$. The profile is symmetric, $a_{n-k}=a_k$, and
$\sum_ka_k=N+4=3.4\,n+4$.

\emph{The form.} Put
\begin{equation}\label{eq:R}
  R_n(t):=(2t+n)\,\frac{\prod_{m=1}^{6}\prod_{c\in C_m}(t+\frac12+c)}{\prod_{k=0}^{n}(t+k)^{a_k}},
  \qquad
  S_n:=\int_{\Z_2}R_n'''\Bigl(t+\frac12-\frac n2\Bigr)\dd t .
\end{equation}
In Pochhammer notation,
\[
  R_n(t)=\frac{(2t+n)\,\mathrm{Num}(t)}{(t)_{n+1}\,(t+\frac n{20})_{\frac9{10}n+1}\,(t+\frac n{10})_{\frac45n+1}\,(t+\frac{3n}{20})_{\frac7{10}n+1}},
\]
where
\begin{multline*}
  \mathrm{Num}(t)=\Bigl(t+\frac12+\frac n8\Bigr)_{\frac34n}\Bigl(t+\frac12+\frac{7n}{40}\Bigr)_{\frac{13}{20}n}^2
  \Bigl(t+\frac12+\frac{9n}{40}\Bigr)_{\frac{11}{20}n}\\
  \times\Bigl(t+\frac12+\frac{11n}{40}\Bigr)_{\frac9{20}n}
  \Bigl(t+\frac12+\frac{13n}{40}\Bigr)_{\frac7{20}n}.
\end{multline*}
The numerator consists of six nested well-poised blocks, all centred at
$-\frac n2$ (in the variable $t$), and the denominator of four nested
Pochhammer symbols, which produce the pole-order ramps.

\begin{lemma}[Basic properties]\label{lem:basic}
\begin{enumerate}
\item[(a)] Every zero of the numerator is a half-integer $-\frac12-c$, except
  the zero $-\frac n2$ of $2t+n$, which lowers the order of the pole at
  $-\frac n2$ from $4$ to $3$. We have the partial-fraction expansion
  \[
    R_n(t)=\sum_{k=0}^{n}\sum_{i=1}^{a_k}\frac{r_{i,k}}{(t+k)^i},\qquad
    r_{i,k}=[\varepsilon^{a_k-i}]G_k(\varepsilon),\quad
    G_k(\varepsilon):=\varepsilon^{a_k}R_n(-k+\varepsilon),
  \]
  where $G_k$ is holomorphic for $|\varepsilon|<1$ (and $r_{4,n/2}=0$).
\item[(b)] $\deg R_n=1+N-\sum_ka_k=-3$; hence
  $c_1:=\sum_kr_{1,k}=\lim_{t\to\infty}tR_n(t)=0$.
\item[(c)] $R_n(-n-t)=-R_n(t)$. Consequently $r_{i,n-k}=(-1)^{i+1}r_{i,k}$,
  and $c_i:=\sum_kr_{i,k}=0$ for even $i$.
\item[(d)] $a_k\le4$, so $c_i=0$ for $i\ge5$. Hence $c_3$ is the only
  possibly non-zero $c_i$.
\end{enumerate}
\end{lemma}

\begin{proof}
(a) and (d) are read off the definitions, and (b) by counting degrees. For
(c), the substitution $t\mapsto-n-t$ maps $2t+n$ to $-(2t+n)$, each factor
$t+\frac12+c$ to $-(t+\frac12+(n-1-c))$, and each $t+k$ to $-(t+n-k)$; since
the blocks and the profile are symmetric,
$R_n(-n-t)=(-1)^{1+N+\sum a_k}R_n(t)$, and $N+\sum_ka_k=6.8\,n+4$ is even.
Comparing the principal parts at $-k$ and $-(n-k)$ gives
$r_{i,n-k}=(-1)^{i+1}r_{i,k}$; for even $i$ the terms of $c_i$ cancel in
pairs, and the term $k=\frac n2$ vanishes.
\end{proof}

Three independent exact implementations confirm $c_1=c_2=c_4=0$ and $c_3\ne0$
at every $n$ computed.

\subsection{The linear form}\label{ss:linear}

Put $d_k:=k-\frac n2$, let $H_k$ be the set of half-integers strictly between
$0$ and $d_k$ (so $H_{n/2}=\emptyset$), and
\begin{equation}\label{eq:X}
  X(k,h):=\sum_{i=1}^{a_k}(i)_4\,r_{i,k}\,h^{-i-4}\qquad(h\in H_k).
\end{equation}

\begin{lemma}[The linear form]\label{lem:L}
We have
\begin{equation}\label{eq:L}
  \begin{gathered}
  S_n=\rho_n+Z_n\,\zeta_2(7),\qquad Z_n=-46080\,c_3,\\
  \rho_n=\sum_{k=0}^{n}\sgn(d_k)\sum_{h\in H_k}X(k,h)\in\Q .
  \end{gathered}
\end{equation}
Moreover $\rho_n=-2\sum_{k<n/2}\sum_{h\in H_k}X(k,h)$, and
\begin{equation}\label{eq:X-res}
  X(k,h)=24\,\Res_{t=-k}\,R_n(t)\,(t_e-t)^{-5},\qquad t_e:=-k+h .
\end{equation}
\end{lemma}

\begin{proof}
The poles of $R_n(t+\frac12-\frac n2)$ are half-integers, which are not in
$\Z_2$, so $S_n$ exists, and
$R_n'''(t)=-\sum_{i,k}(i)_3r_{i,k}(t+k)^{-i-3}$. Apply
Lemma~\ref{lem:translation} to $g(t)=(t+\frac12)^{-s}$, with
$g'(l)=-s(l+\frac12)^{-s-1}$ and the shift $d_k$: for $d_k\ge1$,
$\int(t+\frac12+d_k)^{-s}\dd t=J_s-s\sum_{l=0}^{d_k-1}(l+\frac12)^{-s-1}$, and
for $d_k\le-1$ it equals $J_s+s\sum_{l=d_k}^{-1}(l+\frac12)^{-s-1}$. In both
cases the correction is $-\sgn(d_k)\,s\sum_{h\in H_k}h^{-s-1}$. With
$s=i+3$ and $(i)_3(i+3)=(i)_4$, summation gives
\[
  S_n=-\sum_i(i)_3\,c_i\,J_{i+3}+\sum_k\sgn(d_k)\sum_{h\in H_k}\sum_i(i)_4\,r_{i,k}\,h^{-i-4}.
\]
By Lemma~\ref{lem:LSZ28}, $-(i)_3c_iJ_{i+3}=-(i)_42^{i+4}c_i\zeta_2(i+4)$.
The terms $i=1,2,4$ vanish by Lemma~\ref{lem:basic} (and
$\zeta_2(6)=\zeta_2(8)=0$ as well), and the term $i=3$ is
$-360\cdot128\,c_3\zeta_2(7)=-46080\,c_3\zeta_2(7)$.

For the symmetric form, write $\sum_{h\in H_k}X(k,h)=\sum_{t\in T_k}P_k''''(t)$,
where $P_k(t)=\sum_ir_{i,k}(t+k)^{-i}$ is the principal part at $-k$ and
$T_k=H_k-k$ is the set of half-integers strictly between $-k$ and $-\frac n2$.
By Lemma~\ref{lem:basic}(c), $P_{n-k}''''(t)=-P_k''''(-n-t)$; the map
$t\mapsto-n-t$ sends $T_{n-k}$ onto $T_k$, and $\sgn(d_{n-k})=-\sgn(d_k)$. So
the terms for $k$ and $n-k$ are equal. Finally, for \eqref{eq:X-res} put
$u=t+k$, so that $t_e-t=h-u$ and $(t_e-t)^{-5}$ is analytic at $u=0$. Only the
principal part of $R_n$ at $-k$ contributes to the residue, and
$24(h-u)^{-5}=\sum_{i\ge1}(i)_4h^{-i-4}u^{i-1}$, so the residue is
$\sum_ir_{i,k}(i)_4h^{-i-4}/24$.
\end{proof}

\begin{remark}\label{rem:L-checks}
Formula \eqref{eq:L} was implemented by five independent programs, which
agree. A zeta-free evaluation of $S_n$ (from Mahler series or from the Taylor
coefficients of the product form) agrees with $\rho_n-60c_3J_6$ modulo
$2^{v_2(S_n)+64}$ for $n\le1600$, where $J_6$ was computed in three
independent ways (Bernoulli numbers, a Mahler series, and a precomputed
table), which agree to at least $17000$ bits. These checks are not used in
the proof.
\end{remark}

% ===========================================================================
\section{The \texorpdfstring{$2$}{2}-adic size (Theorem V)}\label{sec:V}

We prove a general bound, of which the case needed here is \eqref{eq:V-cand}.
It extends the $2$-adic estimates of \cite[\S4]{Lai25} and \cite{LSZ26} to
shifted blocks and variable pole orders.

\begin{namedthm}{Theorem V}
Let $n\ge1$, $j\ge0$, $\delta\in\{0,1\}$ and $m\in\Z$. Let
$[\mathrm{lo}_\nu,\mathrm{hi}_\nu]$ ($\nu=1,\dots,B$) be integer intervals of
lengths $L_\nu\ge0$, put $x_\nu:=\mathrm{hi}_\nu+1-m$, and let
$a_k\ge0$ ($0\le k\le n$) be integers. Put
\[
  R(t):=(2t+n)^\delta\,\frac{\prod_\nu\prod_{c=\mathrm{lo}_\nu}^{\mathrm{hi}_\nu}(t+\frac12+c)}{\prod_k(t+k)^{a_k}},
  \qquad S:=\int_{\Z_2}R^{(j)}\Bigl(t+\frac12-m\Bigr)\dd t,
\]
$N:=\sum_\nu L_\nu$, $\sigma:=\sum_\nu s_2(L_\nu)$ and
$\lambda:=\lfloor\log_2\max(1,\max_\nu L_\nu)\rfloor+1$. Then
\[
  v_2(S)\ge N+\sum_ka_k+j-\sigma-(j+1)\lambda .
\]
\end{namedthm}

For $R_n$ we have $B=6$, $j=3$, $m=\frac n2$, $x_\nu=(\frac12-\eta_\nu)n\ge0$,
$\max_\nu L_\nu=\frac34n$ and $\sigma\le6\lambda$, so
\begin{equation}\label{eq:V-cand}
  v_2(S_n)\ge6.8\,n+7-\sigma_n-4\lambda_n\ge6.8\,n+7-10\lambda_n,\qquad
  \lambda_n=\Bigl\lfloor\log_2\frac{3n}4\Bigr\rfloor+1 .
\end{equation}

\begin{proof}
\emph{Factorization.} Since
$t+\frac12-m+\frac12+c=t+1+c-m$ and
$t+\frac12-m+k=\frac12(2t+2k+1-2m)$,
\begin{gather*}
  R\Bigl(t+\frac12-m\Bigr)=2^{\sum_ka_k}\,E(t)\prod_\nu L_\nu!\binom{t+x_\nu}{L_\nu},\\
  E(t):=(2t+1+n-2m)^\delta\prod_k(2t+2k+1-2m)^{-a_k}.
\end{gather*}
Every factor $2t+b$ with $b$ odd is a $2$-adic unit on $\Z_2$, and its
inverse is a power series in $t$ with coefficients in $\Z_2$ tending to $0$.
Hence $E_\beta:=E^{(\beta)}/(\beta!\,2^\beta)$ is $\Z_2$-valued with
$\Delta(E_\beta)\ge0$, by (D2) and (D3).

\emph{Terms.} By Leibniz' rule and (D5), $R^{(j)}(t+\frac12-m)$ is a sum of
terms $K\cdot\Phi(t)$ with
$\Phi=E_\beta\prod_\nu\binom{t+x_\nu}{L_\nu-r_\nu}$, where $\beta+\sum_\nu\mu_\nu=j$,
$\mu_\nu\le r_\nu\le L_\nu$, and
\[
  K=2^{\sum a_k}\,\frac{j!}{\prod_\nu\mu_\nu!}\,2^\beta\prod_\nu L_\nu!\,c_{\mu_\nu,r_\nu}.
\]
By Legendre's formula and (D5),
$v_2(K)\ge\sum_ka_k+\beta+(N-\sigma)-(j-\beta)(\lambda-1)$. By (D2)--(D4),
$\Delta(\Phi)\ge-(\lambda-1)$, so $v_2(\int\Phi)\ge-\lambda$ by (D1). Hence
$v_2(K\int\Phi)\ge N+\sum a_k+j-\sigma-(j+1)\lambda+\beta\lambda$, which
proves the theorem.
\end{proof}

\begin{remark}\label{rem:V-sharp}
For $n=480$, $960$ and $2000$ the bound \eqref{eq:V-cand}, with the exact
$\sigma_n$, gives $3214$, $6474$ and $13534$, while the true values of
$v_2(S_n)$ are $3237$, $6499$ and $13560$; the slack is $18$ to $27$ bits for
all $n=40,\dots,2000$. The Lean proof obtains $v_2(S_n)\ge6.8n+4-\sigma_n-4(\lambda_n-1)$,
one more than \eqref{eq:V-cand}.
\end{remark}

% ===========================================================================
\section{The odd denominators (Theorems C, K and A)}\label{sec:den}

\subsection{Set-up}\label{ss:den-setup}

Let $D_n$ be the least odd positive integer with
$D_n\rho_n,\,D_nZ_n\in\Z[\frac12]$. For an odd prime $q$ put
\[
  e_q(n):=v_q(D_n)=\max\bigl(0,-v_q(\rho_n),-v_q(Z_n)\bigr).
\]
Since $46080=2^{10}\cdot3^2\cdot5$, we have $v_q(Z_n)\ge v_q(c_3)$ for all odd
$q$, so $e_q(n)\le\max(0,-v_q(\rho_n),-v_q(c_3))$.

Fix an odd prime $q$ and put $x:=n/q$. A pair $(k,h)$ with $h\in H_k$ is
\emph{bad} if $q\mid2h$ (note that $2h$ is an odd integer) and \emph{good}
otherwise. Write $\rho_n=G+B$, where $G$ and $B$ are the sums in
\eqref{eq:L} over the good and the bad pairs.

We begin with the primes above $n$, where the centred evaluation pays off.

\begin{namedthm}{Theorem C}
If $q>n$ is an odd prime, then $v_q(\rho_n)\ge0$ and $v_q(c_3)\ge0$; hence
$q\nmid D_n$.
\end{namedthm}

\begin{proof}
All pole differences satisfy $|k'-k|\le n<q$, and the values $h\in H_k$
satisfy $|2h|\le n-1<q$, so these are $q$-units; the factor $n-2k$ is a
$q$-unit or, for $k=\frac n2$, is replaced by $2\varepsilon$, and there we
work with $G_k(\varepsilon)/(2\varepsilon)$. The numerator values
$2(c-k)+1$ may be divisible by $q$. Write
$G_k(\varepsilon)=G_k(0)\prod_\beta(1+\varepsilon/\beta)^{e_\beta}$, where
$\beta$ runs over the numerator values $c-k+\frac12$ (with $e_\beta\ge0$) and
over the pole differences $k'-k$ (with $e_\beta=-a_{k'}$), with the obvious
change at $k=\frac n2$. Since $G_k(0)$ contains the factor $\beta^{e_\beta}$
for every numerator value, $G_k(0)\binom{e_\beta}{j_\beta}\beta^{-j_\beta}$ is
$q$-integral for $j_\beta\le e_\beta$, and it vanishes for $j_\beta>e_\beta$
(compare Lemma~\ref{lem:top}). Hence every $r_{i,k}$ is $q$-integral, so is
every $X(k,h)$ ($h$ is a $q$-unit), and so are $\rho_n$ and $c_3$.
\end{proof}

Numerically, no prime in $(n,2n+2]$ divides $D_n$ for any $n\le40000$ tested.

\subsection{Residue discs and local models}\label{ss:discs}

All zeros and poles of $R_n$ lie in $\frac12\Z\subset\Z_q$.

\emph{Discs.} For $k_0\in\{0,\dots,q-1\}$ let $\xi=\xi(k_0)$ be the class of
$-k_0$, with disc $D_\xi:=-k_0+q\Z_q$ and centre $\hat\xi:=-k_0$. For a
rational function $F$ with poles in $\frac12\Z$ put
$\oint_{D_\xi}F:=\sum_{w\in D_\xi}\Res_{t=w}F$. Every pole $-k$ lies in
exactly one disc, and the evaluation point $t_e=-k+h$ of a pair lies in the
disc of $-k$ if and only if the pair is bad. Hence
$B=\sum_\xi\beta_\xi$, where
$\beta_\xi:=\sum\sgn(d_k)X(k,h)$ over the bad pairs with $k\equiv k_0\pmod q$.

\emph{Items.} Write $R_n(t)=2\prod_w(t-w)^{e_w}$. The \emph{items} $w$ are the
poles $-k$ (with $e_w=-a_k$), the zeros $-\frac12-c$ (with $e_w=\mu(c)$), and
the point $-\frac n2$ of the factor $2t+n$ (with $e_w=+1$), which we keep
separate from the pole at $-\frac n2$.

\emph{Local coordinate.} Put $v:=(t-\hat\xi)/q$. An item $w\in D_\xi$ has the
local point $v_w:=(w-\hat\xi)/q\in\frac12\Z$. Define
\[
  RR_\xi(v):=\prod_{w\in D_\xi}(v-v_w)^{e_w},\qquad
  E_\xi:=\sum_{w\in D_\xi}e_w,\qquad
  \hat U_\xi(t):=2\prod_{w\notin D_\xi}(t-w)^{e_w}.
\]
Then $R_n(t)=q^{E_\xi}\,\hat U_\xi(t)\,RR_\xi(v)$ on $D_\xi$.

The \emph{local model} of $\xi$ is the following data: the poles $v=-j$ for
$k=k_0+jq\le n$, of order $a_k$; the zeros $v=-s/2$, $s$ odd, with
$c=k_0+(sq-1)/2\in[0,n)$, of multiplicity $\mu(c)$; the $\delta$-zero
$v=(2k_0-n)/(2q)$, present when $q\mid n-2k_0$, which sits on the pole
$-\frac n2$ and lowers its net order to $3$; and the bad pairs, that is, a
pole $v_k=-j$ with $d_k\ne0$ and an evaluation point
$v_e=v_k+\sgn(d_k)\sigma/2$ with $\sigma$ odd and $\sigma q<2|d_k|$.
Equivalently, $v_e$ is a half-integer strictly between $v_k$ and
$v_c:=(2k_0-n)/(2q)$.

\begin{lemma}[Outside factor]\label{lem:outside}
$\hat U_\xi(\hat\xi+qv)=\sum_{m\ge0}F_mq^mv^m$ with $F_m\in\Z_q$ and
$F_0\in\Z_q^\times$. The series converges for $|v|_q\le1$ and equals the
rational function there.
\end{lemma}

\begin{proof}
For $w\notin D_\xi$, $u_w:=\hat\xi-w$ is a $q$-unit, and
$(u_w+qv)^{e}=u_w^{e}\sum_r\binom er(qv/u_w)^r$ with $\binom er\in\Z$ for
every $e\in\Z$. A finite product of such series has the stated form, and
$\hat U_\xi$ has no pole on $|v|_q\le1$.
\end{proof}

\begin{lemma}[Sizes and $q$-integrality]\label{lem:sizes}
Every local point (pole, zero, $\delta$-point, evaluation point) lies in
$[-x-\frac12,\frac12]$. Consequently every rational number built below from a
local model (the Taylor and Laurent coefficients of $RR_\xi$, of
$(v_e-v)^{-5}$ and of $\prod(1-v_wy)^{e_w}$, polynomial combinations of
these, and the kernel constants of Theorem~\ref{thm:collective}) has a
denominator composed of $2$, of integers $2|v-v'|\le2x+2$, and of odd
$\sigma<2x$. All of these numbers are $q$-integral once $q>2x+2$, in particular
once $q^2>2n+3q$.
\end{lemma}

\begin{proof}
Poles satisfy $v=-j\in[-x,0]$. For zeros, $c\ge0$ and $k_0\le q-1$ force
$s\ge-1$, and $c<n$ forces $s<2x+1/q$. Evaluation points lie between a pole
and $v_c\le\frac12$. Expanding $RR_\xi(v)(v_e-v)^{-5}$ around $v_k$ only
divides by the differences $v_k-v_w$ and $v_e-v_k$.
\end{proof}

\subsection{The class lemma}\label{ss:class}

For a class $\xi$ let $s$ be a local variable and $y:=1/s$. Put
\[
  A_\xi(s):=\sum_{\text{bad pairs }(k,e)\text{ of }\xi}\sgn(d_k)\,
  \mathrm{PP}_{v_k}\bigl[RR_\xi(v)(v_e-v)^{-5}\bigr](s),
\]
where $\mathrm{PP}_{v_k}$ denotes the principal part at $v_k$, and
\[
  r_\xi(y):=\prod_{w\in D_\xi}(1-v_wy)^{e_w}=y^{E_\xi}RR_\xi(1/y)\in\Q[[y]],\qquad r_\xi(0)=1 .
\]
For a \emph{kernel polynomial} $\kappa(y)=c_0+\sum_{j\ge1}c_jy^j\in\Q[y]$ put
$K(v):=\kappa(1/v)$ and
\[
  D_{\xi,\kappa}(y):=y^{E_\xi}A_\xi(1/y)-r_\xi(y)\,\kappa(y).
\]

\begin{lemma}[Class lemma]\label{lem:class}
Let $q\ge5$ with $q>2x+2$, and let the $c_j$ be $q$-integral. Then
\begin{equation}\label{eq:class}
  \beta_\xi-24q^{-5}\oint_{D_\xi}R_n(t)\,K\Bigl(\frac{t-\hat\xi}q\Bigr)\dd t
  =24q^{E_\xi-4}\sum_{m\ge0}F_mq^m\,[y^{m+E_\xi+1}]D_{\xi,\kappa}(y).
\end{equation}
Consequently the left-hand side has $q$-adic valuation at least
\[
  \Lev_\kappa(\xi):=-5+\min\{p\ge E_\xi+1:[y^p]D_{\xi,\kappa}\ne0\}
\]
($=+\infty$ if there is no such $p$), and always
$\Lev_\kappa(\xi)\ge E_\xi-4$.
\end{lemma}

\begin{proof}
(i) \emph{The bad part.} By \eqref{eq:X-res}, a bad pair has
$X(k,h)=24\Res_{t=-k}R_n(t)(t_e-t)^{-5}$. Substituting $t=\hat\xi+qv$, so that
$\dd t=q\dd v$ and $t_e-t=q(v_e-v)$, gives
$X=24q^{E-4}\Res_{v=v_k}\bigl[\hat U(\hat\xi+qv)RR(v)(v_e-v)^{-5}\bigr]$.
Insert Lemma~\ref{lem:outside}; the re-expansion at $v_k$ converges in $\Q_q$
because $|v_k|_q\le1$. The residue at $v_k$ only sees the principal part
$P_{k,e}$ of $RR(v)(v_e-v)^{-5}$, which is a rational function with its only
pole at $v_k$ and vanishing at $\infty$; hence
$\Res_{v_k}[v^mP_{k,e}]=-\Res_\infty[v^mP_{k,e}]=[y^{m+1}]P_{k,e}(1/y)$.
Summing with signs,
$\beta_\xi=24q^{E-4}\sum_mF_mq^m[y^{m+E+1}]\bigl(y^EA_\xi(1/y)\bigr)$.

(ii) \emph{The subtracted part.} The poles of $R_n(t)K(v)$ in $D_\xi$ are the
local poles, together with $v=0$ if some $c_j\ne0$; these are all the finite
poles of the rational function $v^mRR(v)K(v)$. Hence
$\oint_{D_\xi}R_nK\dd t=q^{E+1}\sum_mF_mq^m\cdot[v^{-1}]_\infty(v^mRR\,K)$,
where $[v^{-1}]_\infty$ is the coefficient in the expansion at $\infty$. Since
$v^mRR(v)K(v)=\sum_p[y^p](r_\xi\kappa)\,v^{m+E-p}$, this coefficient is
$[y^{m+E+1}](r_\xi\kappa)$.

Multiply (ii) by $24q^{-5}$ and subtract. All coefficients are $q$-integral
(Lemma~\ref{lem:sizes} and $F_m\in\Z_q$), and the $m$-th term has valuation at
least $E-4+m$.
\end{proof}

Three kinds of kernels are used: $\kappa=0$, which gives the levels
$\Lev_0(\xi)$ (a class without bad pairs has $\beta_\xi=0$); $\kappa=1$,
which gives $\Lev_1(\xi)$; and $\kappa=1+c\,y^j$ with $j\ge0$, which gives
$\Lev_\kappa(\xi)$. In the computations the levels are evaluated in exact
rational arithmetic with series truncated at order $12$: the coefficients are
exact below the truncation index, and if no non-zero coefficient is found
there, the truncation index is used, which is still a lower bound.

\subsection{The kernel identity and the far part}\label{ss:kernel}

\begin{lemma}[Kernel identity]\label{lem:kernel}
Let $\Psi(t):=c_0+\sum_{j=1}^{J}c_jq^j\sum_{i=0}^{q-1}(t+i)^{-j}$. Then
$\sum_\xi\oint_{D_\xi}R_n\Psi\dd t=0$.
\end{lemma}

\begin{proof}
$R_n\Psi$ is a rational function whose finite poles are integers, each lying
in exactly one disc. So the left-hand side is the sum of all finite residues,
which equals $-\Res_\infty(R_n\Psi)=0$ because
$\deg(R_n\Psi)\le\deg R_n=-3\le-2$.
\end{proof}

On $D_\xi$ the summand $i=k_0$ of $\Psi$ equals $c_jq^j(t-\hat\xi)^{-j}=c_jv^{-j}$.
Hence $\Psi=K(v)+\sum_{j\ge1}c_jq^jg_{j,\xi}(t)$ with
$g_{j,\xi}(t):=\sum_{i\ne k_0}(t+i)^{-j}$, which is analytic on $D_\xi$ since
$0<|i-k_0|<q$.

\begin{lemma}[Far part]\label{lem:farpart}
Let $q>2x+2$, $j\ge1$, $c_j$ $q$-integral, and $q\ge j+1-E_\xi$. Then
\[
  v_q\Bigl(24q^{-5}\,c_jq^j\oint_{D_\xi}R_n\,g_{j,\xi}\dd t\Bigr)\ge j-4 .
\]
\end{lemma}

\begin{proof}
We have $g_{j,\xi}(\hat\xi+qv)=\sum_{r\ge0}\binom{-j}rq^rv^rP_{j+r}$ with
$P_s:=\sum_{i\ne k_0}(i-k_0)^{-s}$. The numbers $i-k_0$ ($i\ne k_0$) form a
complete system of non-zero residues modulo $q$, so $P_s\equiv0\pmod q$ unless
$(q-1)\mid s$. Write $\hat U\cdot g=\sum_mH_mq^mv^m$. For $m\le q-2-j$,
every $P_{j+r}$ with $r\le m$ vanishes modulo $q$, so $q\mid H_m$; for larger
$m$, $v_q(H_mq^m)\ge m\ge q-1-j$. As in part~(ii) of the proof of
Lemma~\ref{lem:class}, $\oint_{D_\xi}R_ng=q^{E+1}\sum_mH_mq^m[y^{m+E+1}]r_\xi$,
and only $m\ge\max(0,-E-1)$ contribute. Each term has valuation at least
$E+1+m+1\ge1$ (first range) or $E+1+q-1-j\ge1$ (second range, by the
hypothesis on $q$). Multiplying by $24q^{j-5}c_j$ gives the claim.
\end{proof}

\begin{theorem}[Collective bound for the bad part]\label{thm:collective}
Let $q\ge29$ with $q>2x+2$. Put
\[
  L_a:=\min_{\xi\text{ exposed}}\Lev_0(\xi),\qquad L_b:=\min_\xi\Lev_1(\xi),
\]
where a class is \emph{exposed} if it has a bad pair. Define $L_\kappa$ as
follows. Put $P_0:=L_b$ and $j:=P_0+5$. If $j\ge0$ and every local type of a
class with $E_\xi+1\le j$ has the same coefficient
$c:=[y^j]D_{\xi,1}\ne0$, set $\kappa:=1+c\,y^j$ and
$L_\kappa:=\min(\min_\xi\Lev_\kappa(\xi),P_0+1)$; otherwise set
$L_\kappa:=P_0$. Then
\[
  v_q(B)\ge\mathfrak B_q:=\max(L_a,L_b,L_\kappa).
\]
\end{theorem}

\begin{proof}
By Lemma~\ref{lem:kernel},
$B=\sum_\xi\bigl[\beta_\xi-24q^{-5}\oint_{D_\xi}R_n\Psi\bigr]$ for every
admissible $\Psi$: $\Psi=0$ gives the bound $L_a$, $\Psi=1$ gives $L_b$, and
the $\kappa$-kernel gives $L_\kappa$ (for $j=0$ it is the constant
$\Psi=1+c$, with no far part; for $j\ge1$ it is $c_0=1$, $c_j=c$). Each
bracket is \eqref{eq:class}, plus, for $j\ge1$, the far part of
Lemma~\ref{lem:farpart}, which costs $j-4=P_0+1$. For every type with
$E+1\le j$ we have $[y^p]D_{\xi,\kappa}=[y^p]D_{\xi,1}=0$ for $E+1\le p<j$
(by the definition of $P_0$) and $[y^j]D_{\xi,\kappa}=c-c=0$, so
$\Lev_\kappa\ge P_0+1$. Types with $E+1>j$ have
$\Lev_\kappa\ge-5+E+1\ge P_0+1$ whatever $\kappa$ is. Lemma~\ref{lem:farpart}
applies, because $E_\xi\ge-10$ (Lemma~\ref{lem:crude}) and $j\le4$ (the step
only matters when $P_0<0$), so $q\ge29\ge j+1-E_\xi$.
\end{proof}

\begin{remark}\label{rem:j-negative}
Only $j\ge0$ is admissible. For $j<0$ the ``kernel'' would be a polynomial in
$v$, the degree of $R_n\Psi$ could be $\ge-1$, and
Lemmas~\ref{lem:kernel}--\ref{lem:farpart} would fail. An early version of the
computation overlooked this and produced bounds below the true exponent at
$19$ primes; the verification step detected it.
\end{remark}

\subsection{The good part and \texorpdfstring{$c_3$}{c3}}\label{ss:good}

\begin{lemma}[Per-pole bound]\label{lem:perpole}
Let $q\ge5$ with $q>2x+2$. Let $-k$ be a pole in the class $\xi$ with net local
order $a$ ($a=a_k$, or $a=3$ at $-\frac n2$), and write
$\tau_k(w):=w^aRR_\xi(v_k+w)=\sum_l\tau_lw^l$ ($\tau_0\ne0$). Put
$\ell_k(i):=\max\{l\le a-i:\tau_l\ne0\}$. Then
$v_q(r_{i,k})\ge E_\xi+a-\ell_k(i)$.
\end{lemma}

\begin{proof}
Expand around $t_k=-k$ as in Lemma~\ref{lem:outside}:
\[
  R_n(t_k+qw)=q^E\,\hat U(t_k+qw)\,w^{-a}\tau_k(w),\qquad
  \hat U(t_k+qw)=\sum_mF'_mq^mw^m,\quad F'_m\in\Z_q .
\]
So
$r_{i,k}=q^{E+i}\sum_{l+m=a-i}F'_mq^m\tau_l$, and every term has valuation at
least $E+a-l$ ($\tau_l$ is $q$-integral by Lemma~\ref{lem:sizes}).
\end{proof}

Let $\mathfrak G_q$ be the minimum of $E+a-\ell(i)$ over the local types
realized by the $q$ classes, their poles and $i\le a$, and let $\mathfrak C_q$
be the minimum of $E+a-\ell(3)$ over the poles with $a\ge3$. For $q\ge11$ the
numbers $(i)_4\in\{24,120,360,840\}$ are $q$-units, and so is $h$ for a good
pair; hence $v_q(G)\ge\mathfrak G_q$ and $v_q(c_3)\ge\mathfrak C_q$.

\subsection{Theorem K}\label{ss:K}

\begin{namedthm}{Theorem K}
Let $q\ge29$ be a prime with $q^2>2n+3q$. Then
\[
  e_q(n)\le f_n(q):=\max\bigl(0,\,-\min(\mathfrak B_q,\mathfrak G_q),\,-\mathfrak C_q\bigr).
\]
\end{namedthm}

\begin{proof}
$q^2>2n+3q$ gives $q>2x+3$, so Lemmas~\ref{lem:sizes}--\ref{lem:perpole} and
Theorem~\ref{thm:collective} apply. Then $\rho_n=G+B$ with
$v_q(B)\ge\mathfrak B_q$ and $v_q(G)\ge\mathfrak G_q$, and
$v_q(c_3)\ge\mathfrak C_q$.
\end{proof}

The number $f_n(q)$ is a finite exact computation on the local models. It was
implemented twice, independently (once directly, once through residue moments
modulo the prime $2^{127}-1$).

\subsection{Crude bounds and the small primes}\label{ss:crude}

\begin{lemma}[Crude level bounds and exposure]\label{lem:crude}
For every $n$, every prime $q\ge5$ and every class:
\begin{enumerate}
\item[(1)] $E_\xi\ge-10$. Hence, if $q\ge29$ and $q^2>2n+3q$, then
  $e_q(n)\le14$.
\item[(2)] If $x>3$, every class is exposed. If $q\ge29$ and $q^2>2n+3q$,
  then $e_q(n)\le\max(0,4-E_{\min}(n,q))$, where $E_{\min}(n,q)$ is the
  minimum of $E_\xi$ over the $q$ classes.
\end{enumerate}
\end{lemma}

\begin{proof}
(1) Block $m$ contributes at least $\lfloor L_m/q\rfloor>L_m/q-1$ zeros to the
class, because its values $2(c-k_0)+1$ are $L_m$ consecutive odd numbers.
Level $\ell$ of the pole profile contributes fewer than $|I_\ell|/q+1$ poles,
where $I_\ell:=\{k:a_k\ge\ell\}$. Since $\sum_mL_m=N$ and
$\sum_\ell|I_\ell|=N+4$, this gives $E>-10-4/q$, so $E\ge-10$. Then
$L_a\ge E-4\ge-14$, $\mathfrak G\ge E+1$ and $\mathfrak C\ge E+3$, and
Theorem~K gives $e_q(n)\le14$.

(2) The pole $k=k_0\le q-1$ of the class has
$|d_{k_0}|/q\ge x/2-1+1/q>\frac12$, so $\sigma=1$ gives a bad pair. The bound
follows from Theorem~K with $L_a\ge E_{\min}-4$, $\mathfrak G\ge E_{\min}+1$
and $\mathfrak C\ge E_{\min}+3$ (exposure is not needed for it: classes
without bad pairs contribute nothing to $B$).
\end{proof}

\begin{lemma}[Small primes]\label{lem:small}
For every prime $q\ge5$ and every $n$, $e_q(n)\le18\lfloor\log_q(2n+1)\rfloor$,
and $e_3(n)\le19\lfloor\log_3(2n+1)\rfloor$.
\end{lemma}

\begin{proof}
Put $\lambda_1:=\lfloor\log_q(2n+1)\rfloor$. For $k\ne\frac n2$ write
$G_k(\varepsilon)=G_k(0)\prod_\beta(1+\varepsilon/\beta)^{e_\beta}$ with
$\beta\in\frac12\Z\setminus\{0\}$ and $|2\beta|\le2n+1$; so
$v_q(\beta)\le\lambda_1$, and Lemma~\ref{lem:top} gives
$v_q([\varepsilon^\lambda]\prod)\ge-\lambda\lambda_1$. For $k=\frac n2$, where
the linear factor is $2\varepsilon$ and $G_k(0)=0$, apply the same argument to
$G_k(\varepsilon)/(2\varepsilon)$; since $H_{n/2}=\emptyset$, only
$r_{3,n/2}$ matters there, and the count below is unchanged.

For $h\in H_k$ we have $|2h|<n$, and
$X(k,h)=24[\varepsilon^{a_k-1}]G_k(\varepsilon)(h-\varepsilon)^{-5}$ with
$[\varepsilon^r](h-\varepsilon)^{-5}=\binom{r+4}4h^{-5-r}$. Hence
$v_q(X(k,h))\ge s_k-8\lambda_1$ and $v_q(r_{3,k})\ge s_k-\lambda_1$, where
$s_k:=v_q(G_k(0))$.

The valuation $s_k$ is at least the numerator count minus the denominator
count. In the numerator, the $L_m$ values of block $m$ are consecutive odd
numbers, so their $q$-part is at least $\sum_e\lfloor L_m/q^e\rfloor$. In the
denominator, for each level $\ell$ with $I_\ell\cap\Z=[\alpha,\beta']$, the
product $\prod_{k'\in[\alpha,\beta'],k'\ne k}|k'-k|$ equals
$(\beta'-\alpha)!/\binom{\beta'-\alpha}{k-\alpha}$ if $k\in[\alpha,\beta']$,
and $|I_\ell|!\binom{\beta'-k}{|I_\ell|}$ (or the mirror image) otherwise; its
$q$-part is at most $v_q(|I_\ell|!)+\lfloor\log_qn\rfloor$ by Kummer's bound.
So
\[
  s_k\ge\sum_{e\ge1}\Bigl[\sum_m\Bigl\lfloor\frac{L_m}{q^e}\Bigr\rfloor-\sum_\ell\Bigl\lfloor\frac{|I_\ell|}{q^e}\Bigr\rfloor\Bigr]-4\lfloor\log_qn\rfloor .
\]
Each bracket is an integer $>-6-4/q^e$, hence $\ge-6$ ($\ge-7$ if $q=3$), and
it vanishes once $q^e>n+1$. Hence $s_k\ge-10\lambda_1$ ($-11\lambda_1$ if
$q=3$), and the lemma follows.
\end{proof}

Consequently $\sum_{q\le Q}e_q(n)\log q\le19\,\pi(Q)\log(2n+1)=O(\sqrt n)$
for $Q:=\max(40,\sqrt{3n})$.

\subsection{The continuum function}\label{ss:continuum}

With $x=n/q$ and $y:=k_0/q\in[0,1)$ the local model of a class is described
by $(x,y)$ as follows.
\begin{itemize}
\item \emph{Poles.} The pole $j$ exists if and only if $y+j\le x$; its order
  is $a(\kappa)$ with $\kappa=(y+j)/x$, by the rule of \S\ref{ss:construction}
  applied to $\min(\kappa,1-\kappa)$.
\item \emph{Zeros.} The zero at $v=-s/2$ has multiplicity
  $\#\{m:\eta_mx<y+s/2<(1-\eta_m)x\}$. The discrete end points are never
  hit, since $2k_0+sq$ is odd while $2\eta_mn$ is even.
\item \emph{Bad pairs.} $d_k/q=y+j-\frac x2$, and $(k,\sigma)$ is a bad pair
  if and only if $\sigma<2|y+j-\frac x2|$; equality never occurs, since
  $|2k-n|$ is even and $\sigma q$ is odd.
\item \emph{Centre class.} This is the class $y\equiv\frac x2\pmod1$, where
  the $\delta$-zero sits on the pole at $v=-\lfloor x/2\rfloor$.
\end{itemize}
As $y$ runs over $[0,1)$ with $x$ fixed, the local model changes only when $y$
crosses a point $\equiv\gamma x+e\pmod1$ with $\gamma\in\{b/40:0\le b\le40\}$
and $e\in\frac12\Z$, or crosses $y\equiv0$. These \emph{events} come from pole
existence ($\gamma=1$), from the cells $b/40$ of the pole profile, from the
block ends $\eta_m$, $1-\eta_m$, and from exposure and the centre
($\gamma=\frac12$). The \emph{candidate set} is
$\mathcal C:=\bigcup_{d=1}^{40}\frac{20}d\Z$.

\begin{lemma}[Piecewise constancy]\label{lem:piecewise}
Two events $\gamma x+e$ and $\gamma'x+e'$ coincide modulo $1$ only if
$(\gamma-\gamma')x\in\frac12\Z$, that is, $x\in\frac{20}d\Z$ with
$d=40|\gamma-\gamma'|$; and an event meets $y\equiv0$ only if
$\gamma x\in\frac12\Z$. Hence the cyclic order of the events is constant on
each open interval of $\R\setminus\mathcal C$. In particular the set
$\mathcal T(x)$ of \emph{continuum types} (the local model in each gap
between consecutive events, together with the centre type) and the
\emph{continuum bound}
\[
  f(x):=\max\bigl(0,-\min(\mathfrak B(\mathcal T(x)),\mathfrak G(\mathcal T(x))),-\mathfrak C(\mathcal T(x))\bigr)
\]
are constant on each such interval.
\end{lemma}

\begin{lemma}[Discrete types are continuum types]\label{lem:discrete}
Let $q>40$ be a prime with $q\nmid n/40$. Then $x=n/q\notin\mathcal C$, every
local type realized by the $q$ classes at $(n,q)$ lies in $\mathcal T(x)$, and
$f_n(q)\le f(x)$. Every prime $q>n/40$ satisfies $q\nmid n/40$.
\end{lemma}

\begin{proof}
(i) If $x=20m'/d$ with $d\le40$, then $2(n/40)d=m'q$; since $q>40\ge d$ and $q$
is odd, $q\mid n/40$.
(ii) By the parity remarks above, a discrete class never sits on a zero event
or an exposure event. A class on a pole-existence event or a cell event takes
that event's one-sided state; since $x\notin\mathcal C$, no other event is at
the same place, so its type is that of an adjacent gap. The class
$k_0\equiv\frac n2$ is the centre type.
(iii) \emph{Monotonicity.} Over fewer types, $L_a$, $L_b$, $\mathfrak G$ and
$\mathfrak C$ can only increase. For $L_\kappa$: if the discrete $P_0$ equals
the continuum $P_0$, some discrete type attains it, so the discrete kernel
step is admissible with the same $(j,c)$, and its levels are among the
continuum ones; if the discrete $P_0$ is larger, then
$\mathfrak B(\text{discrete})\ge P_0(\text{discrete})\ge P_0+1\ge L_\kappa(\text{continuum})$.
\end{proof}

\begin{lemma}[Periodicity]\label{lem:periodic}
$E(x+40,y)=E(x,y)$ for all $(x,y)$, and $\mathcal C+40=\mathcal C$. Hence the
residue level $f^0(x):=\max(0,4-\min_{\mathcal T(x)}E)$ is $40$-periodic, and
$f(x)\le f^0(x)$ for $x>3$.
\end{lemma}

\begin{proof}
Let $I_\ell^c:=\{\kappa\in[0,1]:a(\kappa)\ge\ell\}$; these are the intervals
$[0,1]$, $[\frac1{20},\frac{19}{20}]$, $[\frac1{10},\frac9{10}]$,
$[\frac3{20},\frac{17}{20}]$, of total length $3.4$. Then
\[
  E=\sum_m\#\Bigl\{s\text{ odd}:y+\frac s2\in(\eta_mx,(1-\eta_m)x)\Bigr\}
  -\sum_{\ell=1}^{4}\#\{j\ge0:y+j\in x\,I_\ell^c\}+\delta .
\]
Under $x\mapsto x+40$ every end point $40\eta_m$, $40(1-\eta_m)$ and
$40\,\partial I^c_\ell$ is an integer, so each count grows by exactly
$40(1-2\eta_m)$ or $40|I^c_\ell|$; since
$\sum_m(1-2\eta_m)=3.4=\sum_\ell|I^c_\ell|$, the changes cancel, and the
$\delta$-point $y\equiv x/2$ moves by $20$. The last claim is
Lemma~\ref{lem:crude}(2).
\end{proof}

\subsection{Theorem A}\label{ss:A}

\begin{namedthm}{Theorem A}
As $n\to\infty$ through $40\Z$,
\[
  \frac1n\log D_n\le E^*+o(1),\qquad
  E^*:=\int_1^{60}f(x)\frac{\dd x}{x^2}+\int_{60}^\infty f^0(x)\frac{\dd x}{x^2}=3.88084310467\ldots
\]
\end{namedthm}

\begin{proof}
We split the odd primes into five ranges.
\begin{enumerate}
\item $q>n$: by Theorem~C, $q\nmid D_n$.
\item $\frac n{60}<q\le n$, $q\nmid\frac n{40}$: for $n\ge n_0$ we have $q>40$
  and $q^2>2n+3q$, so $e_q(n)\le f(n/q)$ by Theorem~K and
  Lemma~\ref{lem:discrete}. The function $f$ is constant on the finitely many
  candidate intervals $(a_i,b_i)$ of $[1,60]$, and $n/q$ is never an end
  point (the end points lie in $\mathcal C$). So by the prime number theorem
  $\sum f(n/q)\log q\le\sum_if_i\bigl(\theta(n/a_i)-\theta(n/b_i)\bigr)
  =n\int_1^{60}f\,\frac{\dd x}{x^2}+o(n)$.
\item $\sqrt{3n}<q\le\frac n{60}$, $q\nmid\frac n{40}$: here
  $q^2>3n\ge2n+3q$, so $e_q(n)\le f^0(n/q)$ by Lemma~\ref{lem:crude}(2)
  and Lemma~\ref{lem:discrete}. The function $f^0\le14$ is piecewise constant
  and $40$-periodic. For any $X_2>60$, the prime number theorem on
  $[60,X_2]$ and Chebyshev's bound for $q\le n/X_2$ give at most
  $n\int_{60}^\infty f^0\frac{\dd x}{x^2}+14.3\,n/X_2+o(n)$.
\item $q\mid\frac n{40}$, $q>\sqrt{3n}$: there are at most $\log_2n$ such
  primes, each with $e_q(n)\le14$ (Lemma~\ref{lem:crude}(1)); they contribute
  $O(\log^2n)$.
\item $q\le\max(40,\sqrt{3n})$: $O(\sqrt n)$, by Lemma~\ref{lem:small}.
\end{enumerate}
Letting $X_2\to\infty$ proves the theorem. The value of $E^*$ is computed in
Table~\ref{tab:Estar}.
\end{proof}

\begin{table}[ht]
\centering
\small
\renewcommand{\arraystretch}{1.2}
\caption{The value of $E^*$. The step function $f$ is computed exactly on
each candidate interval. The periodic tail is computed from the exact values
$f^0_i$ of $f^0$ on the $980$ intervals $(a_i,b_i)$ of the period $[60,100)$ as
$\sum_if^0_i(\psi(b_i/40)-\psi(a_i/40))/40$, with the digamma function
$\psi$, $\psi(z)-\psi(w)=\sum_{j\ge0}(\frac1{j+w}-\frac1{j+z})$.}\label{tab:Estar}
\begin{tabular}{llll}
\hline
piece & value & exact value & computation\\
\hline
$\int_1^4f\,\dd x/x^2$ & $2.55$ & $51/20$ & $77$ intervals\\
$\int_4^{20}f\,\dd x/x^2$ & $1.03686695249\ldots$ & $2989163/2882880$ & $392$ intervals\\
$\int_{20}^{60}f\,\dd x/x^2$ & $0.19208902168\ldots$ & see below & $980$ intervals\\
$\int_{60}^\infty f^0\,\dd x/x^2$ & $0.10188713049\ldots$ & -- & one period, digamma sum\\
\hline
$E^*$ & $3.88084310467\ldots$ & & \\
\hline
\end{tabular}
\end{table}

The exact value of the third piece is
\[
  \int_{20}^{60}f\,\frac{\dd x}{x^2}=\frac{5263504492683945763}{27401381122591476000}.
\]

\begin{remark}\label{rem:A-details}
(1) The function $f^0$ has mean $6.0753$ and maximum $10$ over one period. The
tail is not delicate: using only $f^0\le10$ beyond $60$ gives
$E\le3.9456$.
(2) On $[1,3.2)$, $f$ is the step function with values $0,1,2,3,4,5,4,2,3,4,5$
and breaks at $\frac{20}{19},\frac{10}9,\frac{20}{17},\frac54,\frac43,
\frac{20}{13},2,\frac{20}9,\frac{40}{17},\frac{20}7$; its integral there
against $\dd x/x^2$ is $181/80$.
(3) The kernel step with $j\ge1$ (Lemma~\ref{lem:farpart}) is used only on
$[2,\frac{20}9)$, with $j=2$ and $c=-\frac54$; the bounds $\mathfrak G$ and
$\mathfrak C$ are never binding on $[1,60]$.
(4) Where the saving comes from: with the same tail, the residue level
$4-E$ alone gives the exponent $5.0818$; the class lemma with $\kappa=0$
gives $4.3593$; adding the identity with $\kappa=1$ gives $3.9843$; the
one-step kernel gives $E^*=3.8808$.
\end{remark}

\begin{remark}[Evidence and a second proof]\label{rem:A-evidence}
The following checks are not used in the proof. The bound of Theorem~K was
compared with the true exponents $e_q(n)$, computed from exact odd
denominators by independent programs, at $11665$ pairs $(n,q)$ for $22$
values of $n$ up to $200000$ (including every prime at $n=20000$), and at
$1431$ further primes for $n=480,\dots,20000$: there was no violation, and
the bound is attained at $11560$ of the $11665$ pairs. Of the $1449$ head
intervals, the true exponent attains $f$ on $775$, which carry $97.3\%$ of
$\int_1^{60}f\dd x/x^2$; so $E^*$ is essentially sharp for this method. The
class lemma was tested as an exact identity (valuation and leading digit
modulo $q$) on $19827$ class instances. A second, independent proof of a
bound of the same kind, which chooses the kernel subtractions by exact
Gaussian elimination and uses a first-order bound beyond $x=40$, gives the
exponent $3.8801844$ (margin $-0.1265142$) and was compared prime by prime
($19063$ comparisons, no violation).
\end{remark}

% ===========================================================================
\section{The archimedean size (Theorem G)}\label{sec:G}

Let $a(z)$ be the continuum pole profile on $[0,1]$: $a(z)=1,2,3$ for
$\min(z,1-z)$ in $[0,\frac1{20})$, $[\frac1{20},\frac1{10})$,
$[\frac1{10},\frac3{20})$, and $a(z)=4$ otherwise. Define the potential
\[
  U(\kappa):=\sum_{m=1}^{6}\int_{\eta_m}^{1-\eta_m}\log|z-\kappa|\dd z-\int_0^1a(z)\log|z-\kappa|\dd z
  \qquad(\kappa\in[0,1]).
\]
With $F(x):=x\log|x|-x$ this is a finite sum $U(\kappa)=\sum_jw_jF(p_j-\kappa)$
of $18$ terms with integer weights $w_j$ and break points
$p_j\in\frac1{40}\Z\cap[0,1]$.

\begin{lemma}[G1, Cauchy's estimate]\label{lem:G1}
We have $|r_{i,k}|\le2^{-i}M_k$, where
\[
  M_k:=\max_{|\varepsilon|=1/2}|R_n(-k+\varepsilon)| .
\]
\end{lemma}

\begin{proof}
$G_k$ is holomorphic on $|\varepsilon|<1$, and Cauchy's estimate on
$|\varepsilon|=\frac12$ \cite{Ahl79} gives
$|[\varepsilon^\lambda]G_k|\le2^{\lambda-a_k}M_k$. Take $\lambda=a_k-i$.
\end{proof}

\begin{lemma}[G2, triangle inequality]\label{lem:G2}
\begin{multline*}
  \log M_k\le\log(|n-2k|+1)+a_k\log2+\sum_m\sum_{c\in C_m}\log\Bigl(\Bigl|c-k+\frac12\Bigr|+\frac12\Bigr)\\
  -\sum_{k'\ne k}a_{k'}\log\Bigl(|k'-k|-\frac12\Bigr).
\end{multline*}
\end{lemma}

\begin{lemma}[G3, sums and integrals]\label{lem:G3}
Let $f(y):=\log|y-k|$ and $C_3:=2(1-\log2)+\pi^2/72+2=2.750783\ldots$
\begin{enumerate}
\item[(1)] For an integer interval $[\mathrm{lo},\mathrm{hi}]$ of length $L$,
  $\sum_{c=\mathrm{lo}}^{\mathrm{hi}}\log(|c-k+\frac12|+\frac12)\le\int_{\mathrm{lo}}^{\mathrm{hi}+1}f+C_3+\log(2L)$.
\item[(2)] $\sum_{k'\ne k}a_{k'}\log(|k'-k|-\frac12)\ge\int_{-1/2}^{n+1/2}\tilde a_nf-8(1+\log n)$,
  where $\tilde a_n:=a_{k'}$ on $[k'-\frac12,k'+\frac12)$.
\end{enumerate}
\end{lemma}

\begin{proof}
(1) Split $\log(|u|+\frac12)=\log|u|+\log(1+1/(2|u|))$, $u=c-k+\frac12$. The
second part is at most $\sum1/|2(c-k)+1|\le2+\log(2L)$. For the first part
use the midpoint rule on the cells $[c,c+1]$: each of the two cells adjacent
to $k$ contributes at most $1-\log2$, and a cell at distance $d\ge1$ at most
$1/(24d^2)$, since $|f''|\le1/d^2$ there.
(2) Use $\log(d-\frac12)\ge\log d-1/d$ and
$\sum_{k'\ne k}a_{k'}/|k'-k|\le8H_n$. On each cell not containing $k$, $f$ is
concave, so its midpoint value is at least its mean. The omitted cell of $k$
contributes $-a_k(1+\log2)\le0$ to the integral.
\end{proof}

\begin{lemma}[G4--G5, scaling and discretization]\label{lem:G45}
Substituting $y=nz$ gives
$\int_{\alpha n}^{\beta n}f=(\beta-\alpha)n\log n+n\int_\alpha^\beta\log|z-k/n|\dd z$,
and similarly for $\tilde a_n$. The discrete profile $\tilde a_n(nz)$ differs
from $a(z)$ only on $8$ intervals (the six interior jumps and the two end
half-cells), each of length at most $1/(2n)$, and there by $1$; hence
$n\bigl|\int(\tilde a_n(nz)-a(z))\log|z-\kappa|\dd z\bigr|\le4(1+1/(2n)+\log4n)$.
\end{lemma}

\begin{proposition}[G6]\label{prop:G6}
For all $0\le k\le n$ and $n\ge40$,
\[
  \log M_k\le n\,U(k/n)+\log(n+1)+14\log n+37.44 .
\]
\end{proposition}

\begin{proof}
Combine Lemmas~\ref{lem:G2}--\ref{lem:G45}. The two $\log n$ terms of the
scaling combine to $(N-\sum_ka_k)\log n=-4\log n$, and the constants add up
to $37.386+2/n\le37.44$ (with $\sum_m\log(2-4\eta_m)=0.5635$ from the
$\log(2L)$ terms).
\end{proof}

\begin{lemma}[G7, certified maximum]\label{lem:G7}
$0.7067022241\le\max_{[0,1]}U\le U^*:=0.70670223$; the maximum is attained at
$\kappa^*=0.1067347196\ldots$ and at $1-\kappa^*$.
\end{lemma}

The proof is a computer-assisted certification, described in
\S\ref{ss:comp-U}.

\begin{namedthm}{Theorem G}
For every $n\in40\Z_{>0}$,
\[
  \log\max(|\rho_n|,|Z_n|)\le U^*n+14\log n+2\log(n+1)+47.5 .
\]
\end{namedthm}

\begin{proof}
By \eqref{eq:L} and Lemma~\ref{lem:G1},
\begin{align*}
  |\rho_n|&\le\sum_{k,i}(i)_42^{-i}M_k\sum_{h\in H_k}|h|^{-i-4}
  \le16\sum_{i=1}^{4}(i)_4(1-2^{-i-4})\zeta(i+4)\sum_kM_k\\
  &\le21514\,(n+1)\max_kM_k,
\end{align*}
using $\sum_h|h|^{-s}\le2^s(1-2^{-s})\zeta(s)$ for distinct half-integers of
one sign. Also $|Z_n|=46080|\sum_kr_{3,k}|\le5760(n+1)\max_kM_k$. By
Proposition~\ref{prop:G6}, Lemma~\ref{lem:G7} and $\log21600+37.44\le47.5$,
the bound follows.
\end{proof}

\begin{remark}\label{rem:G-true}
The true growth of $\max(|\rho_n|,|Z_n|)$ is about $e^{0.684n}$, smaller than
$e^{U^*n}$: the largest residues sit in a stretch of pole order $3$ where
their signs alternate, and they cancel (measured: $0.6837$ at $n=96000$). A
saddle-point argument could lower $U^*$ by about $0.023$; this is not needed.
\end{remark}

% ===========================================================================
\section{Nonvanishing (Theorem N)}\label{sec:N}

Let
\[
  I:=\{40m:\ 38m-1\text{ prime}\}=\{40,120,160,240,400,600,640,720,840,960,\dots\}.
\]
By Dirichlet's theorem $I$ is infinite, and for $n\in I$ the number
$\ell:=\frac{19}{20}n-1=38m-1\ge37$ is prime.

\begin{namedthm}{Theorem N}
Let $n\in40\Z_{>0}$ and suppose that $\ell:=\frac{19}{20}n-1$ is prime. Then
\[
  v_\ell(\rho_n)=-1,\qquad v_\ell(Z_n)\ge0,\qquad
  \ell\rho_n\equiv\frac{6864}{7}\cdot\frac{r_{1,0}}{\ell^4}\pmod\ell,
\]
where $u_0:=r_{1,0}/\ell^4$ is an $\ell$-unit.
\end{namedthm}

\begin{proof}
Throughout, $\ell\ge37$, $2n<3\ell$ and $\ell^2>2n$.

\emph{Size facts.} (S1) $|k'-k|\le n<2\ell$, so $\ell\mid k'-k$ if and only
if $k'-k=\pm\ell$. (S2) The numerator values $2(c-k)+1$ are odd with absolute
value at most $2n-1<3\ell$, so $\ell$ divides one only if it equals $\pm\ell$,
that is, $c=k+\frac{\ell-1}2$ or $c=k-\frac{\ell+1}2$, and then with
valuation $1$. (S3) $n-2k$ is even and $|n-2k|<2\ell$, so $\ell\nmid n-2k$
unless $k=\frac n2$. (S4) For $h\in H_k$, $|2h|\le n-1<2\ell$, so $\ell\mid2h$
if and only if $2h=\pm\ell$, which forces $|k-\frac n2|\ge\frac{19}{40}n$,
that is, $k\le\frac n{40}$ or $k\ge\frac{39}{40}n$. (S5) Every block contains
the window $[\frac{19}{40}n-1,\frac{21}{40}n]$, and the numbers $(i)_4$ and
$46080$ are $\ell$-units.

\emph{N1 (local data).} Write $s_k:=v_\ell(G_k(0))$.
(a) \emph{End poles} $0\le k\le\frac n{20}+1$. The value
$c=k+\frac{19}{40}n-1$ lies in the window, hence in all six blocks, and gives
$2(c-k)+1=\ell$; the other branch $c=k-\frac{19}{40}n<0$ lies in no block. The
pole partner $k+\ell$ has order $2$ for $k\in\{0,1\}$ and order $1$ for
$k\ge2$. So $s_0=s_1=6-2=4$ and $s_k=5$ for $2\le k\le\frac n{20}+1$. By
Lemma~\ref{lem:top}, $v_\ell(r_{i,k})\ge s_k-(a_k-i)\ge3$. The mirror poles
obey the same bounds by Lemma~\ref{lem:basic}(c).
(b) \emph{Middle poles} $\frac n{20}+1<k<\ell$. Here the pole partners
$k\pm\ell$ lie outside $[0,n]$, so only numerator factors ($e_\beta\ge0$) are
divisible by $\ell$, and Lemma~\ref{lem:top} gives $v_\ell(r_{i,k})\ge0$ (at
$k=\frac n2$ apply it to $G_k(\varepsilon)/(2\varepsilon)$).
(c) \emph{Consequence.} Every $c_i$ is $\ell$-integral, so $v_\ell(Z_n)\ge0$.

\emph{N2 (two terms).} Use $\rho_n=-2\sum_{k<n/2}\sum_{h\in H_k}X(k,h)$. By N1,
the terms with $\ell\nmid2h$ are $\ell$-integral. For $k<\frac n2$ the
remaining terms have $2h=-\ell$, which requires $k\le\frac n{40}$ by (S4);
conversely $h=-\frac\ell2\in H_k$ for all $k\le\frac n{40}$. There $a_k=1$,
and $X_k:=X(k,-\frac\ell2)=24r_{1,k}(-\frac\ell2)^{-5}=-768\,r_{1,k}/\ell^5$
has valuation $s_k-5$, which is $-1$ for $k\in\{0,1\}$ and $\ge0$ otherwise.
Hence $\rho_n\equiv-2(X_0+X_1)\pmod{\Z_{(\ell)}}$, that is,
$\ell\rho_n\equiv1536\,(r_{1,0}+r_{1,1})/\ell^4\pmod\ell$.

\emph{N3 (the ratio).} Since $r_{1,k}=G_k(0)=(n-2k)\,\mathrm{Num}(-k)/\prod_{k'\ne k}(k'-k)^{a_{k'}}$
for $k\in\{0,1\}$, where $\mathrm{Num}$ is the product of the six blocks, the
numerator products telescope blockwise, and the ratio of the pole products
only involves the jumps of the profile:
\[
  \Theta:=\frac{r_{1,1}}{r_{1,0}}
  =\frac{n-2}{n}\prod_{m=1}^{6}\frac{2\eta_mn-1}{2(1-\eta_m)n-1}\cdot
  \frac{-\,n\cdot\frac{17}{20}n\cdot\frac9{10}n\cdot\frac{19}{20}n}
  {(\frac n{20}-1)(\frac n{10}-1)(\frac3{20}n-1)} .
\]
Since $\frac{19}{20}n=\ell+1$, we have $n\equiv\frac{20}{19}\pmod\ell$, and
then $\frac{n-2}n\equiv-\frac9{10}$, the block product is
$\equiv(-\frac78)(-\frac67)^2(-\frac56)(-\frac45)(-\frac34)=\frac9{28}$, and the
pole factor is $\equiv\frac54$. Hence
$\Theta\equiv\theta_0:=-\frac{81}{224}$ and
$1+\theta_0=\frac{143}{224}=\frac{11\cdot13}{2^5\cdot7}$. All factors involve
only the primes $2,3,5,7,11,13,17,19<37\le\ell$, so the reduction is
legitimate.

\emph{Conclusion.} $u_0=r_{1,0}/\ell^4$ is an $\ell$-unit because $s_0=4$, and
$\ell\rho_n\equiv1536(1+\theta_0)u_0=\frac{6864}7u_0\not\equiv0\pmod\ell$,
since $\frac{6864}7=\frac{2^4\cdot3\cdot11\cdot13}7$ is an $\ell$-unit. Hence
$v_\ell(\rho_n)=-1$.
\end{proof}

\begin{remark}\label{rem:N-checks}
The constants were recomputed exactly by a separate program. The leading
residues $\ell\rho_n\bmod\ell$ at the ten $n\in I$ with $n\le960$ are $24$,
$64$, $54$, $26$, $101$, $510$, $271$, $224$, $704$, $434$, as predicted;
Theorem~N was also checked $\ell$-adically at all $334$ values $n\in I$ with
$n\le60000$.
\end{remark}

% ===========================================================================
\section{Proof of the Main Theorem}\label{sec:main}

Suppose that $\zeta_2(7)=u/b$ with $u\in\Z$ and $b\in\Z_{\ge1}$. For $n\in I$
choose $e_n\ge0$ with $A_n:=2^{e_n}D_n\rho_n\in\Z$ and
$B_n:=2^{e_n}D_nZ_n\in\Z$, and put $N_n:=bA_n+uB_n\in\Z$. By
Lemma~\ref{lem:L}, $N_n=b\,2^{e_n}D_n\,S_n$ in $\Q_2$.

\emph{Step 1 ($N_n\ne0$).} Let $n\in I$ with $\ell:=\frac{19}{20}n-1>b$. By
Theorem~N, $v_\ell(bA_n)=v_\ell(D_n)-1$, and $v_\ell(uB_n)\ge v_\ell(D_n)$ (or
$uB_n=0$). So $v_\ell(N_n)=v_\ell(D_n)-1<\infty$ and $N_n\ne0$. (This is
condition~(2) of \cite[Lemma~2.2]{LS25} with $\ell(n)=\frac{19}{20}n-1$.)

\emph{Step 2 (Liouville).} A non-zero integer $N$ satisfies
$1\le|N|\cdot|N|_2$. Here $|N_n|\le(b+|u|)\,2^{e_n}D_n\max(|\rho_n|,|Z_n|)$
and $|N_n|_2\le2^{-e_n}2^{-v_2(S_n)}$, since $D_n$ is odd. Hence
\begin{equation}\label{eq:liouville}
  1\le(b+|u|)\,D_n\max(|\rho_n|,|Z_n|)\,2^{-v_2(S_n)}\qquad
  \text{for all }n\in I\text{ with }\tfrac{19}{20}n-1>b .
\end{equation}

\emph{Step 3.} By Theorem~A (valid for all $n\in40\Z$, hence along $I$),
Theorem~G and \eqref{eq:V-cand}, the logarithm of the right-hand side of
\eqref{eq:liouville} is at most
\[
  \log(b+|u|)+(E^*+U^*-6.8\log2)\,n+o(n)=-0.1258555\ldots\cdot n+o(n)\longrightarrow-\infty .
\]
Since $I$ is infinite, this contradicts \eqref{eq:liouville}. Hence
$\zeta_2(7)$ is irrational. \qed

\medskip
Only the asymptotic inequality $E^*+U^*<6.8\log2$ is used: given $E^*$, any
growth exponent below $6.8\log2-E^*=0.8325577$ suffices, and given $U^*$, any
denominator exponent below $6.8\log2-U^*=4.0066986$. The $o(n)$ of Theorem~A
is ineffective, and we claim no irrationality measure. Table~\ref{tab:margin}
lists the margins obtained with the various bounds.

\begin{table}[ht]
\centering
\small
\renewcommand{\arraystretch}{1.2}
\caption{Margins $E+g-6.8\log2$ per $n$, where $6.8\log2=4.7134008$ and
$U^*=0.70670223$. The last two rows are numerical and not used in the
proof.}\label{tab:margin}
\begin{tabular}{P{1.3cm}P{5.6cm}P{2.0cm}P{2.3cm}}
\hline
growth $g$ & denominator exponent $E$ & margin & status\\
\hline
$U^*$ & $E^*=3.88084310$ (Theorem~A) & $-0.1258555$ & proved\\
$U^*$ & $3.88018442$ (second proof, Remark~\ref{rem:A-evidence}) & $-0.1265142$ & proved\\
$U^*$ & $3.9240892$ (kernel level on $[1,20]$, residue level on $[20,140]$, $10/140$ beyond) & $-0.0826094$ & proved; Lean uses $g=0.71$, $E=3.99$\\
$U^*$ & $3.9308$ (without Lemma~\ref{lem:farpart}) & $-0.0759$ & proved\\
$U^*$ & $3.9843$ (class lemma with $\kappa=0,1$ only) & $-0.0224$ & proved\\
$\approx0.684$ & $\approx3.87\pm0.02$ (true asymptotic value) & $\approx-0.16$ & numerical\\
true & true $D_n$, $480\le n\le6000$ & $-0.23$ to $-0.18$ & numerical\\
\hline
\end{tabular}
\end{table}

\begin{remark}[Finite $n$]\label{rem:finite}
For $n=480$, $640$, $960$, $1440$, $2000$ and $6000$, the quantity
$\frac1n\log\bigl(D_n\max(|\rho_n|,|Z_n|)2^{-v_2(S_n)}\bigr)$ computed from the
exact data is $-0.2296$, $-0.2161$, $-0.2226$, $-0.1802$, $-0.1829$ and
$-0.1768$. Finite $n$ is far from the asymptotic regime for the prime sums:
the large-prime part of $\frac1n\log D_n$ grows from about $3.3$ ($n=480$) to
$3.70$ ($n=6000$), towards $E^*$, while the small-prime part decays like
$n^{-1/2}$. These numbers are not used in the proof.
\end{remark}

% ===========================================================================
\section{The computer-assisted steps}\label{sec:computer}

All computations of the informal proof were done in exact rational arithmetic
(Python \texttt{fractions}) or in interval arithmetic with outward rounding
(\texttt{mpmath.iv} \cite{mpmath}), and every computer-assisted step was
redone by an independent implementation written by a different agent.

\subsection{The maximum of \texorpdfstring{$U$}{U} (Lemma~\ref{lem:G7})}\label{ss:comp-U}

Each term $w_jF(p_j-\kappa)$ is monotone in $\kappa$ on $[-1,1]$, and convex
or concave on each interval between consecutive break points (since
$F''(z)=1/z$). Three certifications were made.
\begin{enumerate}
\item A monotone box bound, a centred first-order bound on boxes without break
  points, and bisection from the grid $\frac1{40}\Z$: $152$ boxes, run twice,
  once in \texttt{mpmath.iv} at $120$ bits with outward rounding and once in
  exact rationals with the program's own enclosures of $\log$ (from the series
  of $\operatorname{artanh}$ with explicit remainders). The negative control
  $0.70670222$ fails, as it must.
\item An independent program bounds the convex terms by chords and the concave
  terms by tangents, with exact rationals and correctly rounded decimal
  logarithms, on $109$ boxes; it finds $\max U=0.7067022242140760\ldots$
\item The Lean proof certifies the weaker bound $\max U\le0.71$ by chords and
  tangents on $42$ boxes with end points in $\frac1{80}\Z$ (worst box bound
  $0.708146$), with rational enclosures of $\log$ checked in Lean
  (\S\ref{sec:formal}).
\end{enumerate}

\subsection{The value of \texorpdfstring{$E^*$}{E*} (Theorem~A)}\label{ss:comp-E}

The function $f$ of Lemma~\ref{lem:piecewise} is evaluated in exact rational
arithmetic on each of the $1449$ candidate intervals of $[1,60]$, from the
continuum types of the interval (Lemma~\ref{lem:piecewise}); for each interval
the program records $f$ and the kernel that certifies it.
The residue level $f^0$ is computed on one period $[20,60]$ ($980$ intervals),
and the tail beyond $60$ is the digamma sum of Table~\ref{tab:Estar}. An
independent implementation, which formulates the levels through residue
moments modulo $2^{127}-1$ and uses its own event set, reproduces all $1449$
interval values exactly, and the three head integrals were re-summed exactly
($51/20$, $2989163/2882880$ and $5263504492683945763/27401381122591476000$).
The formal proof re-checks, in the Lean kernel, a table of $469$ kernel-level
rows on $[1,20)$ and $980$ residue-level rows on $[20,60)$ (extended by
periodicity to $[20,140)$), and the resulting bound $3.9240891610\ldots\le3.99$
(\S\ref{sec:formal}).

\subsection{Other numerical inputs}\label{ss:comp-other}

The constants of Theorem~N ($\theta_0=-\frac{81}{224}$ and
$1536(1+\theta_0)=\frac{6864}7$), of Theorem~G ($21513.3\ldots$, $5760$ and
the constant of Proposition~\ref{prop:G6}) and of Table~\ref{tab:margin} are
hand computations that were rechecked in exact arithmetic.

\subsection{Code and data}\label{ss:code}

The exact reference engines of the research (the computation of $r_{i,k}$,
$\rho_n$, $Z_n$ and of the levels of Section~\ref{sec:den}), the interval
data of $f$ on $[1,20]$ and the interval data of $f^0$ on one period (the
first, second and fourth pieces of Table~\ref{tab:Estar}) are included,
together with an exact-arithmetic mirror of the Lean definitions and an
independent audit script, in the directory \texttt{python/} of the
formalization repository (\S\ref{ss:formal-howto}); the interval data of $f$
on $[20,60]$ and the other certification programs of this section are not.
The formal proof does not depend on any of these programs.

% ===========================================================================
\section{Formal verification}\label{sec:formal}

The Main Theorem is formalized in Lean~4 \cite{MU21} with Mathlib
\cite{Mathlib20}. The development extends the formalization \cite{DevPair26}
of the statement that at least one of $\zeta_2(7)$, $\zeta_2(9)$ is
irrational: it keeps that formalization's trusted definitions, its shared
proofs (the convergence of $J_s$, the translation formula, Lai's criterion,
the $\Delta$-calculus) and its proof of the prime number theorem, and adds
the namespace \texttt{Zeta2.Z7} (directory \texttt{Zeta2Lean/Z7/}) for the
present proof. Lean names are quoted verbatim; in displayed Lean code the
Unicode symbols of the source are typeset as mathematical symbols.

\subsection{The statement}\label{ss:formal-statement}

The main theorem (file \texttt{Zeta2Lean/Z7/Main.lean}) is
\begin{alltt}\small
theorem Zeta2.Z7.zeta2_7_irrational :
    (\(\forall\) s : \(\mathbb{N}\), HasVolkenborn (halfPow s) (J s)) \(\wedge\) \(\neg\) \(\exists\) q : \(\mathbb{Q}\), zeta2 7 = q
\end{alltt}
Here \texttt{q} is coerced from $\Q$ to Mathlib's field
$\Q$\texttt{\_[2]} of $2$-adic numbers, so the second conjunct says that
\texttt{zeta2~7} is irrational in the sense of \S\ref{ss:result}. The first
conjunct says that for every $s$ the Riemann sums of $(x+\frac12)^{-s}$
converge in $\Q_2$ to \texttt{J~s}; it is there because \texttt{J~s} is
defined with Mathlib's \texttt{limUnder}, which returns an unspecified value
when the limit does not exist.

\subsection{The trusted definitions}\label{ss:formal-defs}

The statement uses six definitions (namespace \texttt{Zeta2}, file
\texttt{Zeta2Lean/Defs.lean}, with the namespaces \texttt{Filter},
\texttt{Topology}, \texttt{Finset} and \texttt{PowerSeries} opened). They are verbatim identical to
those of the registered formalization \cite{DevPair26,PalPair26}:
\begin{alltt}\small
def volkenbornSum (f : \(\mathbb{N}\) \(\to\) \(\mathbb{Q}\)_[2]) (N : \(\mathbb{N}\)) : \(\mathbb{Q}\)_[2] :=
  ((2 : \(\mathbb{Q}\)_[2]) ^ N)\(\sp{-1}\) * \(\sum\) x \(\in\) range (2 ^ N), f x

def HasVolkenborn (f : \(\mathbb{N}\) \(\to\) \(\mathbb{Q}\)_[2]) (I : \(\mathbb{Q}\)_[2]) : Prop :=
  Tendsto (volkenbornSum f) atTop (\(\mathcal{N}\) I)

def volkenborn (f : \(\mathbb{N}\) \(\to\) \(\mathbb{Q}\)_[2]) : \(\mathbb{Q}\)_[2] :=
  limUnder atTop (volkenbornSum f)

def halfPow (s x : \(\mathbb{N}\)) : \(\mathbb{Q}\)_[2] :=
  ((x : \(\mathbb{Q}\)_[2]) + 2\(\sp{-1}\))\(\sp{-1}\) ^ s

def J (s : \(\mathbb{N}\)) : \(\mathbb{Q}\)_[2] :=
  volkenborn (halfPow s)

def zeta2 (s : \(\mathbb{N}\)) : \(\mathbb{Q}\)_[2] :=
  J (s - 1) / (((s : \(\mathbb{Q}\)_[2]) - 1) * 2 ^ s)
\end{alltt}
Here $\mathcal N\,I$ is the neighbourhood filter of $I$, so
\texttt{HasVolkenborn~f~I} says that the Riemann sums \eqref{eq:riemann}
converge to $I$. Hence \texttt{J~s} is the integral $J_s$ of \eqref{eq:Js},
and \texttt{zeta2~7}$\,=J_6/768$. All other definitions (the rational
function, $\rho_n$, $Z_n$, the local models, the certificate tables, the
rates) are internal to the proof: if one of them differed from the text, the
formal theorem would still hold as stated.

\subsection{What is cited}\label{ss:formal-cited}

The identification \texttt{zeta2~7}$\,=J_6/768=\zeta_2(7)=L_2(7,\omega^{-6})$
is Lemma~2.8 of \cite{LSZ26} (Lemma~\ref{lem:LSZ28}). It is cited and not
formalized, since Mathlib does not yet contain the Kubota--Leopoldt
$L$-functions (a Lean~3 formalization of $p$-adic $L$-functions exists
\cite{Nar23}, but it is not part of Mathlib). It is needed only to read the
formal theorem as a statement about the Kubota--Leopoldt value, not in its
proof. Since $768$ is a non-zero rational number, given this identification
the formal theorem is equivalent to the Main Theorem, and to the analogous
statement for every normalization that differs by a non-zero rational factor.

\subsection{Structure of the formal proof}\label{ss:formal-structure}

\sloppy
The formalization follows a blueprint. Each ingredient of the informal proof
is a Lean proposition \texttt{Stmt\_X} (file \texttt{Zeta2Lean/Z7/Statements.lean}),
proved in its own file of \texttt{Zeta2Lean/Z7/Proofs/} with the statements it
depends on as hypotheses. The file \texttt{Zeta2Lean/Z7/Assembly.lean} derives
the main theorem, for any rates $g,\delta$ with $g+\delta<6.8\log2$, from the
top-level statements, and \texttt{Main.lean} connects all proofs.
Table~\ref{tab:formal} gives the correspondence. The formal proof differs from
the text in the following points.
\begin{enumerate}
\item[(1)] Sign and indexing: $n=40u$, and the Lean objects \texttt{rho},
  \texttt{Zc} are $-\rho_n$, $-Z_n$ (the integrand is
  $-R_n'''(t+\frac12-\frac n2)$).
\item[(2)] Rates: the growth rate is $g=0.71\ge U^*$ and the denominator rate
  is $\delta=3.99$, with the margin $0.71+3.99=4.70<6.8\log2$
  (\texttt{margin7}, from Mathlib's bound $\log2>0.6931471803$).
\item[(3)] Denominators: the formal Theorem~A uses the kernel-level bound on
  $[1,20)$ ($469$ rows), the residue level on $[20,140)$ ($3\times980$ rows,
  by periodicity), the bound $e_q(n)\le10$ for $n/q\ge140$ (that is,
  $E_\xi\ge-6$), Lemma~\ref{lem:crude}(1) for $q\mid u$, and
  Lemma~\ref{lem:small} for $q\le1000$ or $q^2\le2n+3q$. The table integral
  is $3.9240891610\ldots$, and the kernel-checked bound is $\le3.99$. The
  certificate statements assume $q>1000$.
\item[(4)] Growth: the residue bound $|r_{i,k}|\le C(u+1)^Be^{nU(k/n)}$ is
  proved without complex analysis and without integrals: a weighted norm on
  truncated power series replaces Cauchy's estimate, and exact telescoping
  sums of $F(z)=z\log|z|-z$ replace Lemmas~\ref{lem:G3}--\ref{lem:G45}, with
  weaker but sufficient polynomial factors. The maximum of $U$ is certified as
  in \S\ref{ss:comp-U}(3).
\item[(5)] Nonvanishing: Dirichlet's theorem is Mathlib's
  \texttt{Nat.frequently\_\allowbreak atTop\_\allowbreak prime\_\allowbreak and\_\allowbreak modEq}, and Lai's criterion is
  used in the form ``$|a_{n,j}|\le B_n$ for all $j$, $B_n|L_n|_2\to0$ and
  $L_n\ne0$ for infinitely many $n$'', proved in \cite{DevPair26}.
\end{enumerate}
In the course of the formalization one small gap in the informal proof was
found and repaired: in Lemma~\ref{lem:small} the factorization of $G_k$ fails
at the centre $k=\frac n2$, where $G_k(0)=0$; the text above already contains
the repair (work with $G_k(\varepsilon)/(2\varepsilon)$).

\begin{table}[ht]
\centering
\footnotesize
\setlength{\tabcolsep}{3pt}
\renewcommand{\arraystretch}{1.2}
\caption{The formalization. Lean statements are in the namespace
\texttt{Zeta2.Z7}, and proof files in \texttt{Zeta2Lean/Z7/Proofs/}.}\label{tab:formal}
\begin{tabular}{P{2.9cm}P{4.4cm}P{4.3cm}}
\hline
this paper & Lean statement & proof files\\
\hline
Lemmas~\ref{lem:basic}, \ref{lem:L} & \texttt{Stmt\_PF7}, \texttt{Stmt\_CoeffVanish7},
  \texttt{Stmt\_TranslationNeg}, \texttt{Stmt\_GenLinearFormZ}, \texttt{Stmt\_L7}
 & \texttt{PartialFractions7}, \texttt{CoeffVanish7}, \texttt{TranslationNeg},
   \texttt{GenLinearFormZ}, \texttt{LinearForm7}\\
Theorem~V & \texttt{Stmt\_Valuation7} & \texttt{IntegrandTaylor7}, \texttt{Valuation7}\\
Theorem~C & \texttt{Stmt\_BigPrimes7} & \texttt{BigPrimes7}\\
Lemmas~\ref{lem:class}, \ref{lem:kernel}, \ref{lem:farpart}, \ref{lem:perpole}
 & \texttt{Stmt\_ClassLemma7}, \texttt{Stmt\_KernelIdentity7},
   \texttt{Stmt\_FarPart7}, \texttt{Stmt\_PerPole7}
 & \texttt{ClassLemma7}, \texttt{LocalExpansion7}, \texttt{KernelIdentity7},
   \texttt{FarPart7}, \texttt{PerPole7}\\
Theorem~K, Lemma~\ref{lem:crude}(1) & \texttt{Stmt\_ThmK7}, \texttt{Stmt\_Emin7}
 & \texttt{ThmK7}, \texttt{Emin7}\\
Lemma~\ref{lem:small} & \texttt{Stmt\_SmallPrimes7} & \texttt{SmallPrimes7}\\
Lemmas~\ref{lem:piecewise}--\ref{lem:periodic}, Table~\ref{tab:Estar}
 & \texttt{Stmt\_HeadCert7}, \texttt{Stmt\_ResCert7}, \texttt{Stmt\_ExpBounds7}
 & \texttt{HeadCert7*}, \texttt{ResCert7*}, \texttt{ExpBounds7};
   table \texttt{Zeta2Lean/Z7/DenTable}\\
Theorem~A & \texttt{Stmt\_Denominators7 delta7} & \texttt{Denominators7}\\
Lemmas~\ref{lem:G1}--\ref{lem:G7}, Theorem~G
 & \texttt{Stmt\_ResidueBound7}, \texttt{Stmt\_MaxU7 g7}, \texttt{Stmt\_Growth7 g7}
 & \texttt{ResidueBound7}, \texttt{MaxU7}, \texttt{Growth7}\\
Theorem~N & \texttt{Stmt\_AuxPrime7}, \texttt{Stmt\_Nonvanishing7}
 & \texttt{AuxPrime7}, \texttt{NonvanishingLS7}\\
Section~\ref{sec:main} & \texttt{main7\_of\_stmts}, \texttt{margin7}
 & \texttt{Zeta2Lean/Z7/Assembly}\\
Lemmas~\ref{lem:translation}, \ref{lem:delta}, \ref{lem:lai}; PNT
 & \texttt{Stmt\_Translation}, \texttt{Stmt\_Delta}, \texttt{Stmt\_DeltaFun},
   \texttt{Stmt\_Criterion}, \texttt{Stmt\_JConv}, \texttt{PNT\_Stmt} (namespace \texttt{Zeta2})
 & \texttt{Zeta2Lean/Proofs/*}, \texttt{Zeta2Lean/Cited/PNT}, shared with \cite{DevPair26}\\
Main Theorem & \texttt{zeta2\_7\_irrational} & \texttt{Zeta2Lean/Z7/Main}\\
\hline
\end{tabular}
\end{table}

\subsection{Trust base and checks}\label{ss:formal-trust}

\begin{itemize}
\item \emph{Lean and Mathlib.} Lean \texttt{v4.35.0-rc2} (toolchain
  \texttt{leanprover/\allowbreak lean4:\allowbreak v4.35.0-rc2}) and Mathlib
  at commit \texttt{065356127b}, as in \cite{DevPair26}.
\item \emph{Axioms.} \texttt{\#print axioms} applied to \texttt{Zeta2.Z7.zeta2\_\allowbreak 7\_\allowbreak irrational}
  lists exactly \texttt{propext}, \texttt{Classical.choice} and
  \texttt{Quot.sound}. The script \texttt{scripts/audit.sh} finds no
  \texttt{sorry}, no \texttt{axiom} declaration and no
  \texttt{native\_decide}, \texttt{implemented\_by} or \texttt{extern}.
\item \emph{Numerical certificates.} The denominator table and its integral
  bound are checked by \texttt{decide +kernel}, that is, by reduction in the
  kernel; the bound $\max U\le0.71$ is proved from rational enclosures of
  $\log$ whose side conditions are checked by \texttt{norm\_num}. The Python
  programs that generated these certificates need not be trusted.
\item \emph{Vendored code.} The prime number theorem is proved from Lean
  files of open Mathlib pull requests by T.~Tao, adapted from the
  PrimeNumberTheoremAnd project \cite{PNTLean} (their pull-request
  descriptions state that AI tools were used in writing or golfing the
  code). They are included with their headers and authors; apart from the
  imports and the module header, their code is identical to that of the pull
  requests. Being Lean proofs checked by the kernel, they add nothing to the
  trust base.
\item \emph{Kernel replay.} Lean's \texttt{leanchecker} replayed all $70$
  modules of the import closure of \texttt{Zeta2Lean.Z7.Main} at the
  development commit of \S\ref{ss:formal-howto}, one module at a
  time, without error (at most about $13$\,GB of memory per module, for the
  certificate data).
\item \emph{Mirror and audit.} An exact-arithmetic Python mirror of the Lean
  definitions checks every statement numerically against the research
  programs, and an independent audit script, which shares no code with the
  mirror, tested the statements on about $700{,}000$ instances, with no
  failure. These checks concern the faithfulness of internal definitions and
  are not part of the trust base.
\end{itemize}

\subsection{How to check}\label{ss:formal-howto}

The formalization is published at
\url{https://github.com/gmDevi/zeta2-7-lean}, commit
\texttt{e26813f04b3e84e136c88c00753567e1b442d5db}. It is the development
described here, repackaged for the Palomar registry: every file uses the Lean
module system, unused files were removed, and the heavy certificates are checked
in small pieces so that every kernel check fits in 32\,GB; no statement changed.
This commit is registered in the Palomar registry of Lean-verified results as
entry PALOMAR-2026-09-30-000018 \cite{PalZ726}.
With \texttt{elan}
installed, the commands
\begin{verbatim}
lake exe cache get
lake build
lake env lean Zeta2Lean/Z7/Main.lean
bash scripts/audit.sh
bash scripts/kernels.sh Zeta2Lean.Z7.Main
\end{verbatim}
fetch the compiled Mathlib, build the project, print the axioms of the main
theorem, run the census of \texttt{sorry} and forbidden constructs, and replay
the import closure of the main theorem with \texttt{leanchecker}, one module at
a time.

% ===========================================================================
\section{AI use and verification}\label{sec:ai}

\subsection{How this work was produced}\label{ss:ai-how}

Apart from the published results it cites and the vendored third-party Lean
files for the prime number theorem \cite{PNTLean}, everything in this paper
was produced by AI agents, namely Anthropic Claude models, under the direction
of the author, in September 2026. The author posed the research question,
directed and supervised the work, and is responsible for this paper. In
detail:
\begin{enumerate}
\item[(1)] \emph{Research.} An orchestrating agent planned the research and
  launched subagents. Research agents explored families of rational functions
  with exact rational arithmetic; the construction of
  Section~\ref{sec:construction} was found by a genetic search over pole-order
  ramps and nested blocks, guided by numerical estimates of the four exponents.
\item[(2)] \emph{Proofs.} Separate agents proved the ingredients: one the
  denominator bound (Theorems~C, K, A), one a second, independent version of
  it, and one the growth, the $2$-adic size and the nonvanishing
  (Theorems~G, V, N).
\item[(3)] \emph{Verification of the informal proof.} Each ingredient was
  re-derived by a separate adversarial verifier agent with its own code. An
  assembly agent combined the ingredients, recomputed every number of the
  margin table, and wrote an AI-written referee report on the assembled
  proof. The verifiers found errors of detail (for instance the inadmissible
  kernels of Remark~\ref{rem:j-negative}, and the interval of the local points
  in Lemma~\ref{lem:sizes}); all were corrected, and none affected the
  conclusions. Their existence means that the text may still contain further
  slips.
\item[(4)] \emph{Formalization.} An architect agent wrote the Lean definitions
  and the statements of all lemmas. Prover agents wrote the proofs against
  these frozen statements and reported problems instead of changing them. An
  auditor agent checked the definitions and statements for faithfulness to the
  informal proof and for soundness (no vacuous statements, no loopholes
  through junk values such as $x/0=0$ or truncated subtraction), with
  independent exact-arithmetic scripts. Integrator agents built the project,
  ran the census of forbidden constructs and the kernel replay.
\item[(5)] \emph{Literature.} The literature search of \S\ref{ss:novelty} was
  carried out by AI agents.
\item[(6)] \emph{This paper} was written by AI agents from the informal proof,
  the Lean repository and the cited papers. The numerical values in it were
  copied from the exact computations of the project.
\end{enumerate}

\subsection{What has been verified, and how}\label{ss:ai-what}

\begin{itemize}
\item \emph{Machine-checked:} the formal statement of
  \S\ref{ss:formal-statement}, in terms of the six definitions of
  \S\ref{ss:formal-defs}, from the three standard axioms, by the Lean kernel,
  with a module-by-module replay by \texttt{leanchecker}.
\item \emph{Cited, not machine-checked:} the identification of the Lean
  definition \texttt{zeta2~7} with the Kubota--Leopoldt value $\zeta_2(7)$
  (Lemma~\ref{lem:LSZ28}, from \cite{LSZ26}).
\item \emph{Checked by AI agents only:} the informal proof as written, in
  particular the sharper constants that the formal proof does not use
  ($U^*=0.70670223$, $E^*=3.8808\ldots$), the numerical remarks, and the
  certification programs of Section~\ref{sec:computer}; these programs were
  written by AI agents and have not been reviewed by a human.
\item \emph{Not done:} no human expert has refereed the mathematics, the Lean
  statements or this paper.
\end{itemize}

\subsection{Suggestions for a human referee}\label{ss:ai-referee}

Since the formal proof covers the Main Theorem completely, a referee who
accepts Lean and Mathlib needs to check (i) that the six definitions of
\S\ref{ss:formal-defs} define the Volkenborn integrals $J_s$ and the number
$J_6/768$; (ii) that the statement of \S\ref{ss:formal-statement} expresses
the Main Theorem for this number; and (iii) Lemma~2.8 of \cite{LSZ26}. For the
informal proof, the new mathematics is the denominator bound of
Section~\ref{sec:den}: the class lemma (Lemma~\ref{lem:class}) is a residue
identity, and the weight of the argument lies in the continuum reduction
(Lemmas~\ref{lem:piecewise}--\ref{lem:discrete}) and in the exact computation
of $f$ on $1449$ intervals. Theorems~V, G and~N are standard in method, and
Theorem~N reduces to the constant $\theta_0=-\frac{81}{224}$.

\subsection*{Acknowledgements}
The formal proof relies on Mathlib, and on the Lean formalization of the
Wiener--Ikehara theorem by T.~Tao and the contributors of the
PrimeNumberTheoremAnd project \cite{PNTLean}.

% ===========================================================================
\raggedright
\begin{thebibliography}{99}
\small
\item[] Theorem and lemma numbers refer to arXiv:2304.00816v1 for
\cite{Lai25}, arXiv:2306.10393v2 for \cite{LS25}, and arXiv:2505.05005 (v1
and v2 have the same numbering) for \cite{LSZ26}.

\bibitem{Ahl79}
L.~V. Ahlfors, \emph{Complex analysis}, 3rd ed., McGraw-Hill, New York, 1979.

\bibitem{Ape79}
R.~Ap\'ery, \emph{Irrationalit\'e de $\zeta(2)$ et $\zeta(3)$}, in: Journ\'ees
Arithm\'etiques (Luminy, 1978), Ast\'erisque \textbf{61}, Soc. Math. France,
Paris, 1979, pp.~11--13.

\bibitem{BR01}
K.~Ball and T.~Rivoal, \emph{Irrationalit\'e d'une infinit\'e de valeurs de la
fonction z\^eta aux entiers impairs}, Invent. Math. \textbf{146} (2001),
193--207.

\bibitem{Beu08}
F.~Beukers, \emph{Irrationality of some $p$-adic $L$-values}, Acta Math. Sin.
(Engl. Ser.) \textbf{24} (2008), no.~4, 663--686; arXiv:math/0603277.

\bibitem{Cal05}
F.~Calegari, \emph{Irrationality of certain $p$-adic periods for small $p$},
Int. Math. Res. Not. \textbf{2005}, no.~20, 1235--1249; arXiv:math/0408214.

\bibitem{CDT20}
F.~Calegari, V.~Dimitrov and Y.~Tang, \emph{$p$-adic Eisenstein series,
arithmetic holonomicity criteria, and irrationality of the $2$-adic period
$\zeta_2(5)$}, talk slides (2020),
\url{https://people.maths.ox.ac.uk/newton/lnts/VDimitrov-LNT.pdf}.

\bibitem{CDT25}
F.~Calegari, V.~Dimitrov and Y.~Tang, \emph{Arithmetic holonomy bounds and
effective Diophantine approximation}, in: Proceedings of the International
Congress of Mathematicians 2026, Vol.~3 (Invited Lectures, Sections 1--4),
SIAM, 2026, pp.~356--376, doi:10.1137/25M1806776; arXiv:2510.04156.

\bibitem{Coh07}
H.~Cohen, \emph{Number Theory, Vol.~II: Analytic and Modern Tools}, Grad. Texts
in Math. \textbf{240}, Springer, New York, 2007.

\bibitem{Col82}
R.~F. Coleman, \emph{Dilogarithms, regulators and $p$-adic $L$-functions},
Invent. Math. \textbf{69} (1982), no.~2, 171--208.

\bibitem{Dev26}
G.~M. Devincenzis, \emph{At least one of the $2$-adic zeta values $\zeta_2(7)$,
$\zeta_2(9)$, $\zeta_2(11)$ is irrational}, Lean~4 formalization (produced with
AI agents), \url{https://github.com/gmDevi/zeta2-7-9-11-lean}, 2026.

\bibitem{DevPair26}
G.~M. Devincenzis, \emph{At least one of the $2$-adic zeta values $\zeta_2(7)$,
$\zeta_2(9)$ is irrational}, Lean~4 formalization (produced with AI agents),
\url{https://github.com/gmDevi/zeta2-7-9-lean}, 2026.

\bibitem{FSZ19}
S.~Fischler, J.~Sprang and W.~Zudilin, \emph{Many odd zeta values are
irrational}, Compos. Math. \textbf{155} (2019), no.~5, 938--952;
arXiv:1803.08905.

\bibitem{Fur04}
H.~Furusho, \emph{$p$-adic multiple zeta values I: $p$-adic multiple
polylogarithms and the $p$-adic KZ equation}, Invent. Math. \textbf{155}
(2004), no.~2, 253--286.

\bibitem{Lai25}
L.~Lai, \emph{On the irrationality of certain $2$-adic zeta values}, Int. J.
Number Theory \textbf{21} (2025), no.~1, 207--235,
doi:10.1142/S1793042125500113; arXiv:2304.00816.

\bibitem{LLS25}
L.~Lai, C.~Lupu and J.~Sprang, \emph{On the irrationality of certain $p$-adic
zeta values}, Res. Math. Sci. \textbf{12} (2025), no.~4, Paper No.~77,
doi:10.1007/s40687-025-00559-x; arXiv:2505.23088.

\bibitem{LS25}
L.~Lai and J.~Sprang, \emph{Many $p$-adic odd zeta values are irrational},
Michigan Math. J., advance publication (2026), doi:10.1307/mmj/20236490;
arXiv:2306.10393.

\bibitem{LSZ26}
L.~Lai, J.~Sprang and W.~Zudilin, \emph{A note on the irrationality of
$\zeta_2(5)$}, Int. Math. Res. Not. IMRN \textbf{2026}, no.~16, rnag180,
doi:10.1093/imrn/rnag180; arXiv:2505.05005.

\bibitem{Lon26a}
C.~D. Long, \emph{Hybrid arithmetic holonomy and twenty-two individual $p$-adic
zeta values}, unrefereed draft,
\url{https://github.com/octonion/p-adic-zeta-irrationality} (commit b46a177,
25 August 2026).

\bibitem{Mathlib20}
The mathlib Community, \emph{The Lean mathematical library}, in: Proceedings
of the 9th ACM SIGPLAN International Conference on Certified Programs and
Proofs (CPP 2020), ACM, New York, 2020, pp.~367--381.

\bibitem{mpmath}
The mpmath development team, \emph{mpmath: a Python library for
arbitrary-precision floating-point arithmetic}, version~1.3.0 (2023),
\url{https://mpmath.org}.

\bibitem{MU21}
L.~de Moura and S.~Ullrich, \emph{The Lean~4 theorem prover and programming
language}, in: Automated Deduction -- CADE~28, Lecture Notes in Comput. Sci.
\textbf{12699}, Springer, Cham, 2021, pp.~625--635.

\bibitem{Nar23}
A.~Narayanan, \emph{Formalization of $p$-adic $L$-functions in Lean~3},
preprint, arXiv:2302.14491 (2023).

\bibitem{Pal26}
Palomar registry of Lean-verified results, entry PALOMAR-2026-09-25-000023,
version~1,
\url{https://palomar-registry.org/entry.html?id=PALOMAR-2026-09-25-000023&version=1}.

\bibitem{PalPair26}
Palomar registry of Lean-verified results, entry PALOMAR-2026-09-25-000026,
version~1,
\url{https://palomar-registry.org/entry.html?id=PALOMAR-2026-09-25-000026&version=1}.

\bibitem{PalZ726}
Palomar registry of Lean-verified results, entry PALOMAR-2026-09-30-000018,
version~1,
\url{https://palomar-registry.org/entry.html?id=PALOMAR-2026-09-30-000018&version=1}.

\bibitem{PNTLean}
Lean~4 proofs of the prime number theorem vendored in the formalizations:
T.~Tao, \emph{Compactly supported functions are dense in Schwartz space},
Mathlib pull request \#43046; T.~Tao, \emph{The Wiener--Ikehara Tauberian
theorem}, Mathlib pull request \#43233; T.~Tao, \emph{The weak prime number
theorem}, Mathlib pull request \#43238 (open as of 25 September 2026; head
commit 78e1b2bbd0 of \#43238, which contains the other two),
\url{https://github.com/leanprover-community/mathlib4/pull/43238}. The code
is adapted from the PrimeNumberTheoremAnd project (J.~F.~A. Balderas,
V.~Beffara, A.~Kontorovich, T.~Tao, R.~Van de Velde, A.~Mellendijk,
A.~Irving and other contributors),
\url{https://github.com/AlexKontorovich/PrimeNumberTheoremAnd}.

\bibitem{RV01}
G.~Rhin and C.~Viola, \emph{The group structure for $\zeta(3)$}, Acta Arith.
\textbf{97} (2001), no.~3, 269--293.

\bibitem{Riv00}
T.~Rivoal, \emph{La fonction z\^eta de Riemann prend une infinit\'e de valeurs
irrationnelles aux entiers impairs}, C. R. Acad. Sci. Paris S\'er. I Math.
\textbf{331} (2000), no.~4, 267--270.

\bibitem{Rob00}
A.~M. Robert, \emph{A Course in $p$-adic Analysis}, Grad. Texts in Math.
\textbf{198}, Springer-Verlag, New York, 2000.

\bibitem{RS62}
J.~B. Rosser and L.~Schoenfeld, \emph{Approximate formulas for some functions
of prime numbers}, Illinois J. Math. \textbf{6} (1962), 64--94.

\bibitem{Spr20}
J.~Sprang, \emph{Linear independence result for $p$-adic $L$-values}, Duke
Math. J. \textbf{169} (2020), no.~18, 3439--3476.

\bibitem{Vol72}
A.~Volkenborn, \emph{Ein $p$-adisches Integral und seine Anwendungen, I},
Manuscripta Math. \textbf{7} (1972), 341--373.

\bibitem{Zud01}
W.~Zudilin, \emph{One of the numbers $\zeta(5)$, $\zeta(7)$, $\zeta(9)$,
$\zeta(11)$ is irrational}, Uspekhi Mat. Nauk \textbf{56} (2001), no.~4,
149--150; English transl., Russian Math. Surveys \textbf{56} (2001), no.~4,
774--776.

\bibitem{Zud04}
W.~Zudilin, \emph{Arithmetic of linear forms involving odd zeta values},
J. Th\'eor. Nombres Bordeaux \textbf{16} (2004), no.~1, 251--291.

\end{thebibliography}

\end{document}
